% macros.tex
% Mathematical notation for graduate notes on toric varieties.
%
% The conventions are chosen to stay reasonably close to
% Cox--Little--Schenck, Toric Varieties.
%
% Load this file after preamble.tex:
%     % preamble.tex
% Shared package imports and environment definitions for the course notes.
% This file is intended to be loaded from main.tex via
%     % preamble.tex
% Shared package imports and environment definitions for the course notes.
% This file is intended to be loaded from main.tex via
%     % preamble.tex
% Shared package imports and environment definitions for the course notes.
% This file is intended to be loaded from main.tex via
%     \input{preamble}

% --------------------------------------------------------------------------
% Encoding, fonts, and typography
% --------------------------------------------------------------------------

\usepackage[T1]{fontenc}
\usepackage[utf8]{inputenc} % Harmless on modern LaTeX; useful on older systems.
\usepackage{lmodern}
\usepackage{microtype}

% --------------------------------------------------------------------------
% Page layout and document structure
% --------------------------------------------------------------------------

\usepackage[
  letterpaper,
  inner=1.25in,
  outer=1.00in,
  top=1.00in,
  bottom=1.10in,
  headheight=15pt
]{geometry}

\usepackage{setspace}
\setstretch{1.05}

\usepackage{emptypage}
\usepackage{fancyhdr}

\pagestyle{fancy}
\fancyhf{}
\fancyhead[LE]{\small\thepage}
\fancyhead[RO]{\small\thepage}
\fancyhead[LO]{\small\nouppercase{\rightmark}}
\fancyhead[RE]{\small\nouppercase{\leftmark}}
\renewcommand{\headrulewidth}{0.4pt}
\renewcommand{\footrulewidth}{0pt}

\fancypagestyle{plain}{
  \fancyhf{}
  \fancyfoot[C]{\small\thepage}
  \renewcommand{\headrulewidth}{0pt}
}

\usepackage{titlesec}

\titleformat{\chapter}[display]
  {\normalfont\bfseries}
  {\Large Chapter \thechapter}
  {0.5ex}
  {\Huge}

\titleformat{name=\chapter,numberless}[display]
  {\normalfont\bfseries}
  {}
  {0pt}
  {\Huge}

\titlespacing*{\chapter}{0pt}{-10pt}{30pt}

% --------------------------------------------------------------------------
% Mathematics
% --------------------------------------------------------------------------

\usepackage{amsmath}
\usepackage{amssymb}
\usepackage{amsthm}
\usepackage{mathtools}
\usepackage{mathrsfs}
\usepackage{bm}
\usepackage{dsfont}

\numberwithin{equation}{chapter}

% --------------------------------------------------------------------------
% Lists, tables, and general utilities
% --------------------------------------------------------------------------

\usepackage{enumitem}
\setlist{
  itemsep=2pt,
  topsep=4pt,
  parsep=0pt
}

\usepackage{booktabs}
\usepackage{array}
\usepackage{longtable}
\usepackage{tabularx}
\usepackage{multicol}
\usepackage{ragged2e}

\usepackage{xcolor}
\usepackage{xparse}
\usepackage{etoolbox}
\usepackage{ifthen}
\usepackage{comment}

% --------------------------------------------------------------------------
% Graphics, diagrams, and commutative diagrams
% --------------------------------------------------------------------------

\usepackage{graphicx}
\graphicspath{{figures/}}

\usepackage{tikz}
\usetikzlibrary{
  arrows.meta,
  calc,
  decorations.markings,
  intersections,
  matrix,
  patterns,
  positioning,
  shapes.geometric
}

\usepackage{tikz-cd}

% Basic TikZ styles for fans, lattices, and polytopes.
\tikzset{
  lattice point/.style={
    circle,
    fill,
    inner sep=1.4pt
  },
  ray/.style={
    -{Latex[length=2.2mm]},
    thick
  },
  fan cone/.style={
    fill=black!8,
    draw=black!55
  },
  polytope/.style={
    thick
  },
  toric divisor/.style={
    thick,
    dashed
  }
}

% --------------------------------------------------------------------------
% Color palette
% --------------------------------------------------------------------------

\definecolor{DictionaryBlue}{HTML}{1F4E79}
\definecolor{SurfaceGreen}{HTML}{3F6B4F}
\definecolor{WarningRed}{HTML}{8B2E2E}
\definecolor{ExampleGold}{HTML}{8A641B}
\definecolor{LightGray}{HTML}{F4F4F4}
\definecolor{PastelYellow}{HTML}{FFF8CC}
\definecolor{PausePurple}{HTML}{6F5A8A}
\definecolor{PredictionTeal}{HTML}{287D7A}
\definecolor{CheatsheetSlate}{HTML}{536878}

% --------------------------------------------------------------------------
% Boxed pedagogical environments
% --------------------------------------------------------------------------

\usepackage[most]{tcolorbox}

\tcbset{
  enhanced,
  breakable,
  boxrule=0.7pt,
  arc=1.5mm,
  left=2mm,
  right=2mm,
  top=1.5mm,
  bottom=1.5mm,
  before skip=8pt,
  after skip=8pt
}

\newtcolorbox{chapteroverview}{
  colback=PastelYellow,
  colframe=ExampleGold!70!black,
  title={Chapter overview},
  fonttitle=\bfseries
}

\newtcolorbox{threedictionaries}{
  colback=DictionaryBlue!4,
  colframe=DictionaryBlue,
  title={The three dictionaries},
  fonttitle=\bfseries
}

\newtcolorbox{surfacelaboratory}{
  colback=SurfaceGreen!5,
  colframe=SurfaceGreen,
  title={Surface laboratory},
  fonttitle=\bfseries
}

\newtcolorbox{warningbox}{
  colback=WarningRed!4,
  colframe=WarningRed,
  title={Warning},
  fonttitle=\bfseries
}

\newtcolorbox{guidingquestion}{
  colback=ExampleGold!5,
  colframe=ExampleGold,
  title={Guiding question},
  fonttitle=\bfseries
}

\newtcolorbox{pauseandthink}{
  colback=PausePurple!5,
  colframe=PausePurple,
  title={Pause and think},
  fonttitle=\bfseries
}

\newtcolorbox{prediction}{
  colback=PredictionTeal!5,
  colframe=PredictionTeal,
  title={Prediction},
  fonttitle=\bfseries
}

\newtcolorbox{dictionarytable}{
  colback=DictionaryBlue!3,
  colframe=DictionaryBlue,
  title={Dictionary summary},
  fonttitle=\bfseries
}

\newtcolorbox{agcheatsheet}{
  colback=CheatsheetSlate!5,
  colframe=CheatsheetSlate,
  title={Algebraic geometry cheatsheet},
  fonttitle=\bfseries
}

% --------------------------------------------------------------------------
% Theorem environments
% --------------------------------------------------------------------------

\theoremstyle{plain}
\newtheorem{theorem}{Theorem}[chapter]
\newtheorem{proposition}[theorem]{Proposition}
\newtheorem{lemma}[theorem]{Lemma}
\newtheorem{corollary}[theorem]{Corollary}

\theoremstyle{definition}
\newtheorem{definition}[theorem]{Definition}
\newtheorem{construction}[theorem]{Construction}
\newtheorem{example}[theorem]{Example}
\newtheorem{exercise}[theorem]{Exercise}
\newtheorem{question}[theorem]{Question}

\theoremstyle{remark}
\newtheorem{remark}[theorem]{Remark}
\newtheorem{observation}[theorem]{Observation}
\newtheorem{convention}[theorem]{Convention}
\newtheorem{notation}[theorem]{Notation}

% Unnumbered variants.
\newtheorem*{theorem*}{Theorem}
\newtheorem*{proposition*}{Proposition}
\newtheorem*{definition*}{Definition}
\newtheorem*{example*}{Example}
\newtheorem*{remark*}{Remark}

% --------------------------------------------------------------------------
% Exercises and solutions
% --------------------------------------------------------------------------

\newcounter{exerciseblock}[chapter]
\renewcommand{\theexerciseblock}{\thechapter.\arabic{exerciseblock}}

\NewDocumentEnvironment{exercises}{O{Exercises}}
{
  \section*{#1}
  \addcontentsline{toc}{section}{#1}
  \begin{enumerate}[label=\textbf{\thechapter.\arabic*.},wide=0pt]
}
{
  \end{enumerate}
}

\newif\ifshowsolutions
\showsolutionsfalse
% To compile instructor solutions, change the previous line to:
% \showsolutionstrue

\NewDocumentEnvironment{solution}{}
{
  \ifshowsolutions
    \par\smallskip
    \noindent\textbf{Solution. }\ignorespaces
  \else
    \expandafter\comment
  \fi
}
{
  \ifshowsolutions
    \par\smallskip
  \else
    \expandafter\endcomment
  \fi
}

% --------------------------------------------------------------------------
% Cross-references
% --------------------------------------------------------------------------

\usepackage{hyperref}

\hypersetup{
  colorlinks=true,
  linkcolor=DictionaryBlue,
  citecolor=SurfaceGreen,
  urlcolor=WarningRed,
  pdfauthor={},
  pdftitle={The Three Dictionaries of Toric Geometry},
  pdfsubject={Notes for Renzo's Fall 2026 course at Colorado State University},
  bookmarksopen=true,
  bookmarksnumbered=true
}

\usepackage[nameinlink,noabbrev,capitalise]{cleveref}

\crefname{theorem}{theorem}{theorems}
\crefname{proposition}{proposition}{propositions}
\crefname{lemma}{lemma}{lemmas}
\crefname{corollary}{corollary}{corollaries}
\crefname{definition}{definition}{definitions}
\crefname{construction}{construction}{constructions}
\crefname{example}{example}{examples}
\crefname{exercise}{exercise}{exercises}
\crefname{question}{question}{questions}
\crefname{remark}{remark}{remarks}
\crefname{observation}{observation}{observations}
\crefname{convention}{convention}{conventions}
\crefname{notation}{notation}{notations}

% --------------------------------------------------------------------------
% Bibliography
% --------------------------------------------------------------------------

\usepackage[
  backend=biber,
  style=alphabetic,
  sorting=nyt,
  maxbibnames=99
]{biblatex}

% The present notes use chapter-level prose references rather than citation
% commands.  This guard keeps the project ready for a future bibliography
% without requiring an empty references.bib file in the current release.
\IfFileExists{references.bib}{\addbibresource{references.bib}}{}

% --------------------------------------------------------------------------
% Index
% --------------------------------------------------------------------------

\usepackage{makeidx}
\makeindex

% --------------------------------------------------------------------------
% Drafting tools
% --------------------------------------------------------------------------

% Margin notes can be used during development:
%     \marginnote{Check this proof.}
\usepackage{marginnote}

\newcommand{\todo}[1]{%
  \marginpar{\raggedright\scriptsize\color{WarningRed}\textbf{TODO:} #1}%
}

% Set to true while drafting to display labels in the margins.
\newif\ifshowlabels
\showlabelsfalse

\ifshowlabels
  \usepackage{showkeys}
\fi

% --------------------------------------------------------------------------
% Miscellaneous formatting
% --------------------------------------------------------------------------

\allowdisplaybreaks
\setcounter{tocdepth}{2}
\setcounter{secnumdepth}{3}

\renewcommand{\qedsymbol}{\(\square\)}

% Prevent isolated lines at page boundaries.
\widowpenalty=10000
\clubpenalty=10000
\displaywidowpenalty=10000


% --------------------------------------------------------------------------
% Encoding, fonts, and typography
% --------------------------------------------------------------------------

\usepackage[T1]{fontenc}
\usepackage[utf8]{inputenc} % Harmless on modern LaTeX; useful on older systems.
\usepackage{lmodern}
\usepackage{microtype}

% --------------------------------------------------------------------------
% Page layout and document structure
% --------------------------------------------------------------------------

\usepackage[
  letterpaper,
  inner=1.25in,
  outer=1.00in,
  top=1.00in,
  bottom=1.10in,
  headheight=15pt
]{geometry}

\usepackage{setspace}
\setstretch{1.05}

\usepackage{emptypage}
\usepackage{fancyhdr}

\pagestyle{fancy}
\fancyhf{}
\fancyhead[LE]{\small\thepage}
\fancyhead[RO]{\small\thepage}
\fancyhead[LO]{\small\nouppercase{\rightmark}}
\fancyhead[RE]{\small\nouppercase{\leftmark}}
\renewcommand{\headrulewidth}{0.4pt}
\renewcommand{\footrulewidth}{0pt}

\fancypagestyle{plain}{
  \fancyhf{}
  \fancyfoot[C]{\small\thepage}
  \renewcommand{\headrulewidth}{0pt}
}

\usepackage{titlesec}

\titleformat{\chapter}[display]
  {\normalfont\bfseries}
  {\Large Chapter \thechapter}
  {0.5ex}
  {\Huge}

\titleformat{name=\chapter,numberless}[display]
  {\normalfont\bfseries}
  {}
  {0pt}
  {\Huge}

\titlespacing*{\chapter}{0pt}{-10pt}{30pt}

% --------------------------------------------------------------------------
% Mathematics
% --------------------------------------------------------------------------

\usepackage{amsmath}
\usepackage{amssymb}
\usepackage{amsthm}
\usepackage{mathtools}
\usepackage{mathrsfs}
\usepackage{bm}
\usepackage{dsfont}

\numberwithin{equation}{chapter}

% --------------------------------------------------------------------------
% Lists, tables, and general utilities
% --------------------------------------------------------------------------

\usepackage{enumitem}
\setlist{
  itemsep=2pt,
  topsep=4pt,
  parsep=0pt
}

\usepackage{booktabs}
\usepackage{array}
\usepackage{longtable}
\usepackage{tabularx}
\usepackage{multicol}
\usepackage{ragged2e}

\usepackage{xcolor}
\usepackage{xparse}
\usepackage{etoolbox}
\usepackage{ifthen}
\usepackage{comment}

% --------------------------------------------------------------------------
% Graphics, diagrams, and commutative diagrams
% --------------------------------------------------------------------------

\usepackage{graphicx}
\graphicspath{{figures/}}

\usepackage{tikz}
\usetikzlibrary{
  arrows.meta,
  calc,
  decorations.markings,
  intersections,
  matrix,
  patterns,
  positioning,
  shapes.geometric
}

\usepackage{tikz-cd}

% Basic TikZ styles for fans, lattices, and polytopes.
\tikzset{
  lattice point/.style={
    circle,
    fill,
    inner sep=1.4pt
  },
  ray/.style={
    -{Latex[length=2.2mm]},
    thick
  },
  fan cone/.style={
    fill=black!8,
    draw=black!55
  },
  polytope/.style={
    thick
  },
  toric divisor/.style={
    thick,
    dashed
  }
}

% --------------------------------------------------------------------------
% Color palette
% --------------------------------------------------------------------------

\definecolor{DictionaryBlue}{HTML}{1F4E79}
\definecolor{SurfaceGreen}{HTML}{3F6B4F}
\definecolor{WarningRed}{HTML}{8B2E2E}
\definecolor{ExampleGold}{HTML}{8A641B}
\definecolor{LightGray}{HTML}{F4F4F4}
\definecolor{PastelYellow}{HTML}{FFF8CC}
\definecolor{PausePurple}{HTML}{6F5A8A}
\definecolor{PredictionTeal}{HTML}{287D7A}
\definecolor{CheatsheetSlate}{HTML}{536878}

% --------------------------------------------------------------------------
% Boxed pedagogical environments
% --------------------------------------------------------------------------

\usepackage[most]{tcolorbox}

\tcbset{
  enhanced,
  breakable,
  boxrule=0.7pt,
  arc=1.5mm,
  left=2mm,
  right=2mm,
  top=1.5mm,
  bottom=1.5mm,
  before skip=8pt,
  after skip=8pt
}

\newtcolorbox{chapteroverview}{
  colback=PastelYellow,
  colframe=ExampleGold!70!black,
  title={Chapter overview},
  fonttitle=\bfseries
}

\newtcolorbox{threedictionaries}{
  colback=DictionaryBlue!4,
  colframe=DictionaryBlue,
  title={The three dictionaries},
  fonttitle=\bfseries
}

\newtcolorbox{surfacelaboratory}{
  colback=SurfaceGreen!5,
  colframe=SurfaceGreen,
  title={Surface laboratory},
  fonttitle=\bfseries
}

\newtcolorbox{warningbox}{
  colback=WarningRed!4,
  colframe=WarningRed,
  title={Warning},
  fonttitle=\bfseries
}

\newtcolorbox{guidingquestion}{
  colback=ExampleGold!5,
  colframe=ExampleGold,
  title={Guiding question},
  fonttitle=\bfseries
}

\newtcolorbox{pauseandthink}{
  colback=PausePurple!5,
  colframe=PausePurple,
  title={Pause and think},
  fonttitle=\bfseries
}

\newtcolorbox{prediction}{
  colback=PredictionTeal!5,
  colframe=PredictionTeal,
  title={Prediction},
  fonttitle=\bfseries
}

\newtcolorbox{dictionarytable}{
  colback=DictionaryBlue!3,
  colframe=DictionaryBlue,
  title={Dictionary summary},
  fonttitle=\bfseries
}

\newtcolorbox{agcheatsheet}{
  colback=CheatsheetSlate!5,
  colframe=CheatsheetSlate,
  title={Algebraic geometry cheatsheet},
  fonttitle=\bfseries
}

% --------------------------------------------------------------------------
% Theorem environments
% --------------------------------------------------------------------------

\theoremstyle{plain}
\newtheorem{theorem}{Theorem}[chapter]
\newtheorem{proposition}[theorem]{Proposition}
\newtheorem{lemma}[theorem]{Lemma}
\newtheorem{corollary}[theorem]{Corollary}

\theoremstyle{definition}
\newtheorem{definition}[theorem]{Definition}
\newtheorem{construction}[theorem]{Construction}
\newtheorem{example}[theorem]{Example}
\newtheorem{exercise}[theorem]{Exercise}
\newtheorem{question}[theorem]{Question}

\theoremstyle{remark}
\newtheorem{remark}[theorem]{Remark}
\newtheorem{observation}[theorem]{Observation}
\newtheorem{convention}[theorem]{Convention}
\newtheorem{notation}[theorem]{Notation}

% Unnumbered variants.
\newtheorem*{theorem*}{Theorem}
\newtheorem*{proposition*}{Proposition}
\newtheorem*{definition*}{Definition}
\newtheorem*{example*}{Example}
\newtheorem*{remark*}{Remark}

% --------------------------------------------------------------------------
% Exercises and solutions
% --------------------------------------------------------------------------

\newcounter{exerciseblock}[chapter]
\renewcommand{\theexerciseblock}{\thechapter.\arabic{exerciseblock}}

\NewDocumentEnvironment{exercises}{O{Exercises}}
{
  \section*{#1}
  \addcontentsline{toc}{section}{#1}
  \begin{enumerate}[label=\textbf{\thechapter.\arabic*.},wide=0pt]
}
{
  \end{enumerate}
}

\newif\ifshowsolutions
\showsolutionsfalse
% To compile instructor solutions, change the previous line to:
% \showsolutionstrue

\NewDocumentEnvironment{solution}{}
{
  \ifshowsolutions
    \par\smallskip
    \noindent\textbf{Solution. }\ignorespaces
  \else
    \expandafter\comment
  \fi
}
{
  \ifshowsolutions
    \par\smallskip
  \else
    \expandafter\endcomment
  \fi
}

% --------------------------------------------------------------------------
% Cross-references
% --------------------------------------------------------------------------

\usepackage{hyperref}

\hypersetup{
  colorlinks=true,
  linkcolor=DictionaryBlue,
  citecolor=SurfaceGreen,
  urlcolor=WarningRed,
  pdfauthor={},
  pdftitle={The Three Dictionaries of Toric Geometry},
  pdfsubject={Notes for Renzo's Fall 2026 course at Colorado State University},
  bookmarksopen=true,
  bookmarksnumbered=true
}

\usepackage[nameinlink,noabbrev,capitalise]{cleveref}

\crefname{theorem}{theorem}{theorems}
\crefname{proposition}{proposition}{propositions}
\crefname{lemma}{lemma}{lemmas}
\crefname{corollary}{corollary}{corollaries}
\crefname{definition}{definition}{definitions}
\crefname{construction}{construction}{constructions}
\crefname{example}{example}{examples}
\crefname{exercise}{exercise}{exercises}
\crefname{question}{question}{questions}
\crefname{remark}{remark}{remarks}
\crefname{observation}{observation}{observations}
\crefname{convention}{convention}{conventions}
\crefname{notation}{notation}{notations}

% --------------------------------------------------------------------------
% Bibliography
% --------------------------------------------------------------------------

\usepackage[
  backend=biber,
  style=alphabetic,
  sorting=nyt,
  maxbibnames=99
]{biblatex}

% The present notes use chapter-level prose references rather than citation
% commands.  This guard keeps the project ready for a future bibliography
% without requiring an empty references.bib file in the current release.
\IfFileExists{references.bib}{\addbibresource{references.bib}}{}

% --------------------------------------------------------------------------
% Index
% --------------------------------------------------------------------------

\usepackage{makeidx}
\makeindex

% --------------------------------------------------------------------------
% Drafting tools
% --------------------------------------------------------------------------

% Margin notes can be used during development:
%     \marginnote{Check this proof.}
\usepackage{marginnote}

\newcommand{\todo}[1]{%
  \marginpar{\raggedright\scriptsize\color{WarningRed}\textbf{TODO:} #1}%
}

% Set to true while drafting to display labels in the margins.
\newif\ifshowlabels
\showlabelsfalse

\ifshowlabels
  \usepackage{showkeys}
\fi

% --------------------------------------------------------------------------
% Miscellaneous formatting
% --------------------------------------------------------------------------

\allowdisplaybreaks
\setcounter{tocdepth}{2}
\setcounter{secnumdepth}{3}

\renewcommand{\qedsymbol}{\(\square\)}

% Prevent isolated lines at page boundaries.
\widowpenalty=10000
\clubpenalty=10000
\displaywidowpenalty=10000


% --------------------------------------------------------------------------
% Encoding, fonts, and typography
% --------------------------------------------------------------------------

\usepackage[T1]{fontenc}
\usepackage[utf8]{inputenc} % Harmless on modern LaTeX; useful on older systems.
\usepackage{lmodern}
\usepackage{microtype}

% --------------------------------------------------------------------------
% Page layout and document structure
% --------------------------------------------------------------------------

\usepackage[
  letterpaper,
  inner=1.25in,
  outer=1.00in,
  top=1.00in,
  bottom=1.10in,
  headheight=15pt
]{geometry}

\usepackage{setspace}
\setstretch{1.05}

\usepackage{emptypage}
\usepackage{fancyhdr}

\pagestyle{fancy}
\fancyhf{}
\fancyhead[LE]{\small\thepage}
\fancyhead[RO]{\small\thepage}
\fancyhead[LO]{\small\nouppercase{\rightmark}}
\fancyhead[RE]{\small\nouppercase{\leftmark}}
\renewcommand{\headrulewidth}{0.4pt}
\renewcommand{\footrulewidth}{0pt}

\fancypagestyle{plain}{
  \fancyhf{}
  \fancyfoot[C]{\small\thepage}
  \renewcommand{\headrulewidth}{0pt}
}

\usepackage{titlesec}

\titleformat{\chapter}[display]
  {\normalfont\bfseries}
  {\Large Chapter \thechapter}
  {0.5ex}
  {\Huge}

\titleformat{name=\chapter,numberless}[display]
  {\normalfont\bfseries}
  {}
  {0pt}
  {\Huge}

\titlespacing*{\chapter}{0pt}{-10pt}{30pt}

% --------------------------------------------------------------------------
% Mathematics
% --------------------------------------------------------------------------

\usepackage{amsmath}
\usepackage{amssymb}
\usepackage{amsthm}
\usepackage{mathtools}
\usepackage{mathrsfs}
\usepackage{bm}
\usepackage{dsfont}

\numberwithin{equation}{chapter}

% --------------------------------------------------------------------------
% Lists, tables, and general utilities
% --------------------------------------------------------------------------

\usepackage{enumitem}
\setlist{
  itemsep=2pt,
  topsep=4pt,
  parsep=0pt
}

\usepackage{booktabs}
\usepackage{array}
\usepackage{longtable}
\usepackage{tabularx}
\usepackage{multicol}
\usepackage{ragged2e}

\usepackage{xcolor}
\usepackage{xparse}
\usepackage{etoolbox}
\usepackage{ifthen}
\usepackage{comment}

% --------------------------------------------------------------------------
% Graphics, diagrams, and commutative diagrams
% --------------------------------------------------------------------------

\usepackage{graphicx}
\graphicspath{{figures/}}

\usepackage{tikz}
\usetikzlibrary{
  arrows.meta,
  calc,
  decorations.markings,
  intersections,
  matrix,
  patterns,
  positioning,
  shapes.geometric
}

\usepackage{tikz-cd}

% Basic TikZ styles for fans, lattices, and polytopes.
\tikzset{
  lattice point/.style={
    circle,
    fill,
    inner sep=1.4pt
  },
  ray/.style={
    -{Latex[length=2.2mm]},
    thick
  },
  fan cone/.style={
    fill=black!8,
    draw=black!55
  },
  polytope/.style={
    thick
  },
  toric divisor/.style={
    thick,
    dashed
  }
}

% --------------------------------------------------------------------------
% Color palette
% --------------------------------------------------------------------------

\definecolor{DictionaryBlue}{HTML}{1F4E79}
\definecolor{SurfaceGreen}{HTML}{3F6B4F}
\definecolor{WarningRed}{HTML}{8B2E2E}
\definecolor{ExampleGold}{HTML}{8A641B}
\definecolor{LightGray}{HTML}{F4F4F4}
\definecolor{PastelYellow}{HTML}{FFF8CC}
\definecolor{PausePurple}{HTML}{6F5A8A}
\definecolor{PredictionTeal}{HTML}{287D7A}
\definecolor{CheatsheetSlate}{HTML}{536878}

% --------------------------------------------------------------------------
% Boxed pedagogical environments
% --------------------------------------------------------------------------

\usepackage[most]{tcolorbox}

\tcbset{
  enhanced,
  breakable,
  boxrule=0.7pt,
  arc=1.5mm,
  left=2mm,
  right=2mm,
  top=1.5mm,
  bottom=1.5mm,
  before skip=8pt,
  after skip=8pt
}

\newtcolorbox{chapteroverview}{
  colback=PastelYellow,
  colframe=ExampleGold!70!black,
  title={Chapter overview},
  fonttitle=\bfseries
}

\newtcolorbox{threedictionaries}{
  colback=DictionaryBlue!4,
  colframe=DictionaryBlue,
  title={The three dictionaries},
  fonttitle=\bfseries
}

\newtcolorbox{surfacelaboratory}{
  colback=SurfaceGreen!5,
  colframe=SurfaceGreen,
  title={Surface laboratory},
  fonttitle=\bfseries
}

\newtcolorbox{warningbox}{
  colback=WarningRed!4,
  colframe=WarningRed,
  title={Warning},
  fonttitle=\bfseries
}

\newtcolorbox{guidingquestion}{
  colback=ExampleGold!5,
  colframe=ExampleGold,
  title={Guiding question},
  fonttitle=\bfseries
}

\newtcolorbox{pauseandthink}{
  colback=PausePurple!5,
  colframe=PausePurple,
  title={Pause and think},
  fonttitle=\bfseries
}

\newtcolorbox{prediction}{
  colback=PredictionTeal!5,
  colframe=PredictionTeal,
  title={Prediction},
  fonttitle=\bfseries
}

\newtcolorbox{dictionarytable}{
  colback=DictionaryBlue!3,
  colframe=DictionaryBlue,
  title={Dictionary summary},
  fonttitle=\bfseries
}

\newtcolorbox{agcheatsheet}{
  colback=CheatsheetSlate!5,
  colframe=CheatsheetSlate,
  title={Algebraic geometry cheatsheet},
  fonttitle=\bfseries
}

% --------------------------------------------------------------------------
% Theorem environments
% --------------------------------------------------------------------------

\theoremstyle{plain}
\newtheorem{theorem}{Theorem}[chapter]
\newtheorem{proposition}[theorem]{Proposition}
\newtheorem{lemma}[theorem]{Lemma}
\newtheorem{corollary}[theorem]{Corollary}

\theoremstyle{definition}
\newtheorem{definition}[theorem]{Definition}
\newtheorem{construction}[theorem]{Construction}
\newtheorem{example}[theorem]{Example}
\newtheorem{exercise}[theorem]{Exercise}
\newtheorem{question}[theorem]{Question}

\theoremstyle{remark}
\newtheorem{remark}[theorem]{Remark}
\newtheorem{observation}[theorem]{Observation}
\newtheorem{convention}[theorem]{Convention}
\newtheorem{notation}[theorem]{Notation}

% Unnumbered variants.
\newtheorem*{theorem*}{Theorem}
\newtheorem*{proposition*}{Proposition}
\newtheorem*{definition*}{Definition}
\newtheorem*{example*}{Example}
\newtheorem*{remark*}{Remark}

% --------------------------------------------------------------------------
% Exercises and solutions
% --------------------------------------------------------------------------

\newcounter{exerciseblock}[chapter]
\renewcommand{\theexerciseblock}{\thechapter.\arabic{exerciseblock}}

\NewDocumentEnvironment{exercises}{O{Exercises}}
{
  \section*{#1}
  \addcontentsline{toc}{section}{#1}
  \begin{enumerate}[label=\textbf{\thechapter.\arabic*.},wide=0pt]
}
{
  \end{enumerate}
}

\newif\ifshowsolutions
\showsolutionsfalse
% To compile instructor solutions, change the previous line to:
% \showsolutionstrue

\NewDocumentEnvironment{solution}{}
{
  \ifshowsolutions
    \par\smallskip
    \noindent\textbf{Solution. }\ignorespaces
  \else
    \expandafter\comment
  \fi
}
{
  \ifshowsolutions
    \par\smallskip
  \else
    \expandafter\endcomment
  \fi
}

% --------------------------------------------------------------------------
% Cross-references
% --------------------------------------------------------------------------

\usepackage{hyperref}

\hypersetup{
  colorlinks=true,
  linkcolor=DictionaryBlue,
  citecolor=SurfaceGreen,
  urlcolor=WarningRed,
  pdfauthor={},
  pdftitle={The Three Dictionaries of Toric Geometry},
  pdfsubject={Notes for Renzo's Fall 2026 course at Colorado State University},
  bookmarksopen=true,
  bookmarksnumbered=true
}

\usepackage[nameinlink,noabbrev,capitalise]{cleveref}

\crefname{theorem}{theorem}{theorems}
\crefname{proposition}{proposition}{propositions}
\crefname{lemma}{lemma}{lemmas}
\crefname{corollary}{corollary}{corollaries}
\crefname{definition}{definition}{definitions}
\crefname{construction}{construction}{constructions}
\crefname{example}{example}{examples}
\crefname{exercise}{exercise}{exercises}
\crefname{question}{question}{questions}
\crefname{remark}{remark}{remarks}
\crefname{observation}{observation}{observations}
\crefname{convention}{convention}{conventions}
\crefname{notation}{notation}{notations}

% --------------------------------------------------------------------------
% Bibliography
% --------------------------------------------------------------------------

\usepackage[
  backend=biber,
  style=alphabetic,
  sorting=nyt,
  maxbibnames=99
]{biblatex}

% The present notes use chapter-level prose references rather than citation
% commands.  This guard keeps the project ready for a future bibliography
% without requiring an empty references.bib file in the current release.
\IfFileExists{references.bib}{\addbibresource{references.bib}}{}

% --------------------------------------------------------------------------
% Index
% --------------------------------------------------------------------------

\usepackage{makeidx}
\makeindex

% --------------------------------------------------------------------------
% Drafting tools
% --------------------------------------------------------------------------

% Margin notes can be used during development:
%     \marginnote{Check this proof.}
\usepackage{marginnote}

\newcommand{\todo}[1]{%
  \marginpar{\raggedright\scriptsize\color{WarningRed}\textbf{TODO:} #1}%
}

% Set to true while drafting to display labels in the margins.
\newif\ifshowlabels
\showlabelsfalse

\ifshowlabels
  \usepackage{showkeys}
\fi

% --------------------------------------------------------------------------
% Miscellaneous formatting
% --------------------------------------------------------------------------

\allowdisplaybreaks
\setcounter{tocdepth}{2}
\setcounter{secnumdepth}{3}

\renewcommand{\qedsymbol}{\(\square\)}

% Prevent isolated lines at page boundaries.
\widowpenalty=10000
\clubpenalty=10000
\displaywidowpenalty=10000

%     % macros.tex
% Mathematical notation for graduate notes on toric varieties.
%
% The conventions are chosen to stay reasonably close to
% Cox--Little--Schenck, Toric Varieties.
%
% Load this file after preamble.tex:
%     % preamble.tex
% Shared package imports and environment definitions for the course notes.
% This file is intended to be loaded from main.tex via
%     % preamble.tex
% Shared package imports and environment definitions for the course notes.
% This file is intended to be loaded from main.tex via
%     \input{preamble}

% --------------------------------------------------------------------------
% Encoding, fonts, and typography
% --------------------------------------------------------------------------

\usepackage[T1]{fontenc}
\usepackage[utf8]{inputenc} % Harmless on modern LaTeX; useful on older systems.
\usepackage{lmodern}
\usepackage{microtype}

% --------------------------------------------------------------------------
% Page layout and document structure
% --------------------------------------------------------------------------

\usepackage[
  letterpaper,
  inner=1.25in,
  outer=1.00in,
  top=1.00in,
  bottom=1.10in,
  headheight=15pt
]{geometry}

\usepackage{setspace}
\setstretch{1.05}

\usepackage{emptypage}
\usepackage{fancyhdr}

\pagestyle{fancy}
\fancyhf{}
\fancyhead[LE]{\small\thepage}
\fancyhead[RO]{\small\thepage}
\fancyhead[LO]{\small\nouppercase{\rightmark}}
\fancyhead[RE]{\small\nouppercase{\leftmark}}
\renewcommand{\headrulewidth}{0.4pt}
\renewcommand{\footrulewidth}{0pt}

\fancypagestyle{plain}{
  \fancyhf{}
  \fancyfoot[C]{\small\thepage}
  \renewcommand{\headrulewidth}{0pt}
}

\usepackage{titlesec}

\titleformat{\chapter}[display]
  {\normalfont\bfseries}
  {\Large Chapter \thechapter}
  {0.5ex}
  {\Huge}

\titleformat{name=\chapter,numberless}[display]
  {\normalfont\bfseries}
  {}
  {0pt}
  {\Huge}

\titlespacing*{\chapter}{0pt}{-10pt}{30pt}

% --------------------------------------------------------------------------
% Mathematics
% --------------------------------------------------------------------------

\usepackage{amsmath}
\usepackage{amssymb}
\usepackage{amsthm}
\usepackage{mathtools}
\usepackage{mathrsfs}
\usepackage{bm}
\usepackage{dsfont}

\numberwithin{equation}{chapter}

% --------------------------------------------------------------------------
% Lists, tables, and general utilities
% --------------------------------------------------------------------------

\usepackage{enumitem}
\setlist{
  itemsep=2pt,
  topsep=4pt,
  parsep=0pt
}

\usepackage{booktabs}
\usepackage{array}
\usepackage{longtable}
\usepackage{tabularx}
\usepackage{multicol}
\usepackage{ragged2e}

\usepackage{xcolor}
\usepackage{xparse}
\usepackage{etoolbox}
\usepackage{ifthen}
\usepackage{comment}

% --------------------------------------------------------------------------
% Graphics, diagrams, and commutative diagrams
% --------------------------------------------------------------------------

\usepackage{graphicx}
\graphicspath{{figures/}}

\usepackage{tikz}
\usetikzlibrary{
  arrows.meta,
  calc,
  decorations.markings,
  intersections,
  matrix,
  patterns,
  positioning,
  shapes.geometric
}

\usepackage{tikz-cd}

% Basic TikZ styles for fans, lattices, and polytopes.
\tikzset{
  lattice point/.style={
    circle,
    fill,
    inner sep=1.4pt
  },
  ray/.style={
    -{Latex[length=2.2mm]},
    thick
  },
  fan cone/.style={
    fill=black!8,
    draw=black!55
  },
  polytope/.style={
    thick
  },
  toric divisor/.style={
    thick,
    dashed
  }
}

% --------------------------------------------------------------------------
% Color palette
% --------------------------------------------------------------------------

\definecolor{DictionaryBlue}{HTML}{1F4E79}
\definecolor{SurfaceGreen}{HTML}{3F6B4F}
\definecolor{WarningRed}{HTML}{8B2E2E}
\definecolor{ExampleGold}{HTML}{8A641B}
\definecolor{LightGray}{HTML}{F4F4F4}
\definecolor{PastelYellow}{HTML}{FFF8CC}
\definecolor{PausePurple}{HTML}{6F5A8A}
\definecolor{PredictionTeal}{HTML}{287D7A}
\definecolor{CheatsheetSlate}{HTML}{536878}

% --------------------------------------------------------------------------
% Boxed pedagogical environments
% --------------------------------------------------------------------------

\usepackage[most]{tcolorbox}

\tcbset{
  enhanced,
  breakable,
  boxrule=0.7pt,
  arc=1.5mm,
  left=2mm,
  right=2mm,
  top=1.5mm,
  bottom=1.5mm,
  before skip=8pt,
  after skip=8pt
}

\newtcolorbox{chapteroverview}{
  colback=PastelYellow,
  colframe=ExampleGold!70!black,
  title={Chapter overview},
  fonttitle=\bfseries
}

\newtcolorbox{threedictionaries}{
  colback=DictionaryBlue!4,
  colframe=DictionaryBlue,
  title={The three dictionaries},
  fonttitle=\bfseries
}

\newtcolorbox{surfacelaboratory}{
  colback=SurfaceGreen!5,
  colframe=SurfaceGreen,
  title={Surface laboratory},
  fonttitle=\bfseries
}

\newtcolorbox{warningbox}{
  colback=WarningRed!4,
  colframe=WarningRed,
  title={Warning},
  fonttitle=\bfseries
}

\newtcolorbox{guidingquestion}{
  colback=ExampleGold!5,
  colframe=ExampleGold,
  title={Guiding question},
  fonttitle=\bfseries
}

\newtcolorbox{pauseandthink}{
  colback=PausePurple!5,
  colframe=PausePurple,
  title={Pause and think},
  fonttitle=\bfseries
}

\newtcolorbox{prediction}{
  colback=PredictionTeal!5,
  colframe=PredictionTeal,
  title={Prediction},
  fonttitle=\bfseries
}

\newtcolorbox{dictionarytable}{
  colback=DictionaryBlue!3,
  colframe=DictionaryBlue,
  title={Dictionary summary},
  fonttitle=\bfseries
}

\newtcolorbox{agcheatsheet}{
  colback=CheatsheetSlate!5,
  colframe=CheatsheetSlate,
  title={Algebraic geometry cheatsheet},
  fonttitle=\bfseries
}

% --------------------------------------------------------------------------
% Theorem environments
% --------------------------------------------------------------------------

\theoremstyle{plain}
\newtheorem{theorem}{Theorem}[chapter]
\newtheorem{proposition}[theorem]{Proposition}
\newtheorem{lemma}[theorem]{Lemma}
\newtheorem{corollary}[theorem]{Corollary}

\theoremstyle{definition}
\newtheorem{definition}[theorem]{Definition}
\newtheorem{construction}[theorem]{Construction}
\newtheorem{example}[theorem]{Example}
\newtheorem{exercise}[theorem]{Exercise}
\newtheorem{question}[theorem]{Question}

\theoremstyle{remark}
\newtheorem{remark}[theorem]{Remark}
\newtheorem{observation}[theorem]{Observation}
\newtheorem{convention}[theorem]{Convention}
\newtheorem{notation}[theorem]{Notation}

% Unnumbered variants.
\newtheorem*{theorem*}{Theorem}
\newtheorem*{proposition*}{Proposition}
\newtheorem*{definition*}{Definition}
\newtheorem*{example*}{Example}
\newtheorem*{remark*}{Remark}

% --------------------------------------------------------------------------
% Exercises and solutions
% --------------------------------------------------------------------------

\newcounter{exerciseblock}[chapter]
\renewcommand{\theexerciseblock}{\thechapter.\arabic{exerciseblock}}

\NewDocumentEnvironment{exercises}{O{Exercises}}
{
  \section*{#1}
  \addcontentsline{toc}{section}{#1}
  \begin{enumerate}[label=\textbf{\thechapter.\arabic*.},wide=0pt]
}
{
  \end{enumerate}
}

\newif\ifshowsolutions
\showsolutionsfalse
% To compile instructor solutions, change the previous line to:
% \showsolutionstrue

\NewDocumentEnvironment{solution}{}
{
  \ifshowsolutions
    \par\smallskip
    \noindent\textbf{Solution. }\ignorespaces
  \else
    \expandafter\comment
  \fi
}
{
  \ifshowsolutions
    \par\smallskip
  \else
    \expandafter\endcomment
  \fi
}

% --------------------------------------------------------------------------
% Cross-references
% --------------------------------------------------------------------------

\usepackage{hyperref}

\hypersetup{
  colorlinks=true,
  linkcolor=DictionaryBlue,
  citecolor=SurfaceGreen,
  urlcolor=WarningRed,
  pdfauthor={},
  pdftitle={The Three Dictionaries of Toric Geometry},
  pdfsubject={Notes for Renzo's Fall 2026 course at Colorado State University},
  bookmarksopen=true,
  bookmarksnumbered=true
}

\usepackage[nameinlink,noabbrev,capitalise]{cleveref}

\crefname{theorem}{theorem}{theorems}
\crefname{proposition}{proposition}{propositions}
\crefname{lemma}{lemma}{lemmas}
\crefname{corollary}{corollary}{corollaries}
\crefname{definition}{definition}{definitions}
\crefname{construction}{construction}{constructions}
\crefname{example}{example}{examples}
\crefname{exercise}{exercise}{exercises}
\crefname{question}{question}{questions}
\crefname{remark}{remark}{remarks}
\crefname{observation}{observation}{observations}
\crefname{convention}{convention}{conventions}
\crefname{notation}{notation}{notations}

% --------------------------------------------------------------------------
% Bibliography
% --------------------------------------------------------------------------

\usepackage[
  backend=biber,
  style=alphabetic,
  sorting=nyt,
  maxbibnames=99
]{biblatex}

% The present notes use chapter-level prose references rather than citation
% commands.  This guard keeps the project ready for a future bibliography
% without requiring an empty references.bib file in the current release.
\IfFileExists{references.bib}{\addbibresource{references.bib}}{}

% --------------------------------------------------------------------------
% Index
% --------------------------------------------------------------------------

\usepackage{makeidx}
\makeindex

% --------------------------------------------------------------------------
% Drafting tools
% --------------------------------------------------------------------------

% Margin notes can be used during development:
%     \marginnote{Check this proof.}
\usepackage{marginnote}

\newcommand{\todo}[1]{%
  \marginpar{\raggedright\scriptsize\color{WarningRed}\textbf{TODO:} #1}%
}

% Set to true while drafting to display labels in the margins.
\newif\ifshowlabels
\showlabelsfalse

\ifshowlabels
  \usepackage{showkeys}
\fi

% --------------------------------------------------------------------------
% Miscellaneous formatting
% --------------------------------------------------------------------------

\allowdisplaybreaks
\setcounter{tocdepth}{2}
\setcounter{secnumdepth}{3}

\renewcommand{\qedsymbol}{\(\square\)}

% Prevent isolated lines at page boundaries.
\widowpenalty=10000
\clubpenalty=10000
\displaywidowpenalty=10000


% --------------------------------------------------------------------------
% Encoding, fonts, and typography
% --------------------------------------------------------------------------

\usepackage[T1]{fontenc}
\usepackage[utf8]{inputenc} % Harmless on modern LaTeX; useful on older systems.
\usepackage{lmodern}
\usepackage{microtype}

% --------------------------------------------------------------------------
% Page layout and document structure
% --------------------------------------------------------------------------

\usepackage[
  letterpaper,
  inner=1.25in,
  outer=1.00in,
  top=1.00in,
  bottom=1.10in,
  headheight=15pt
]{geometry}

\usepackage{setspace}
\setstretch{1.05}

\usepackage{emptypage}
\usepackage{fancyhdr}

\pagestyle{fancy}
\fancyhf{}
\fancyhead[LE]{\small\thepage}
\fancyhead[RO]{\small\thepage}
\fancyhead[LO]{\small\nouppercase{\rightmark}}
\fancyhead[RE]{\small\nouppercase{\leftmark}}
\renewcommand{\headrulewidth}{0.4pt}
\renewcommand{\footrulewidth}{0pt}

\fancypagestyle{plain}{
  \fancyhf{}
  \fancyfoot[C]{\small\thepage}
  \renewcommand{\headrulewidth}{0pt}
}

\usepackage{titlesec}

\titleformat{\chapter}[display]
  {\normalfont\bfseries}
  {\Large Chapter \thechapter}
  {0.5ex}
  {\Huge}

\titleformat{name=\chapter,numberless}[display]
  {\normalfont\bfseries}
  {}
  {0pt}
  {\Huge}

\titlespacing*{\chapter}{0pt}{-10pt}{30pt}

% --------------------------------------------------------------------------
% Mathematics
% --------------------------------------------------------------------------

\usepackage{amsmath}
\usepackage{amssymb}
\usepackage{amsthm}
\usepackage{mathtools}
\usepackage{mathrsfs}
\usepackage{bm}
\usepackage{dsfont}

\numberwithin{equation}{chapter}

% --------------------------------------------------------------------------
% Lists, tables, and general utilities
% --------------------------------------------------------------------------

\usepackage{enumitem}
\setlist{
  itemsep=2pt,
  topsep=4pt,
  parsep=0pt
}

\usepackage{booktabs}
\usepackage{array}
\usepackage{longtable}
\usepackage{tabularx}
\usepackage{multicol}
\usepackage{ragged2e}

\usepackage{xcolor}
\usepackage{xparse}
\usepackage{etoolbox}
\usepackage{ifthen}
\usepackage{comment}

% --------------------------------------------------------------------------
% Graphics, diagrams, and commutative diagrams
% --------------------------------------------------------------------------

\usepackage{graphicx}
\graphicspath{{figures/}}

\usepackage{tikz}
\usetikzlibrary{
  arrows.meta,
  calc,
  decorations.markings,
  intersections,
  matrix,
  patterns,
  positioning,
  shapes.geometric
}

\usepackage{tikz-cd}

% Basic TikZ styles for fans, lattices, and polytopes.
\tikzset{
  lattice point/.style={
    circle,
    fill,
    inner sep=1.4pt
  },
  ray/.style={
    -{Latex[length=2.2mm]},
    thick
  },
  fan cone/.style={
    fill=black!8,
    draw=black!55
  },
  polytope/.style={
    thick
  },
  toric divisor/.style={
    thick,
    dashed
  }
}

% --------------------------------------------------------------------------
% Color palette
% --------------------------------------------------------------------------

\definecolor{DictionaryBlue}{HTML}{1F4E79}
\definecolor{SurfaceGreen}{HTML}{3F6B4F}
\definecolor{WarningRed}{HTML}{8B2E2E}
\definecolor{ExampleGold}{HTML}{8A641B}
\definecolor{LightGray}{HTML}{F4F4F4}
\definecolor{PastelYellow}{HTML}{FFF8CC}
\definecolor{PausePurple}{HTML}{6F5A8A}
\definecolor{PredictionTeal}{HTML}{287D7A}
\definecolor{CheatsheetSlate}{HTML}{536878}

% --------------------------------------------------------------------------
% Boxed pedagogical environments
% --------------------------------------------------------------------------

\usepackage[most]{tcolorbox}

\tcbset{
  enhanced,
  breakable,
  boxrule=0.7pt,
  arc=1.5mm,
  left=2mm,
  right=2mm,
  top=1.5mm,
  bottom=1.5mm,
  before skip=8pt,
  after skip=8pt
}

\newtcolorbox{chapteroverview}{
  colback=PastelYellow,
  colframe=ExampleGold!70!black,
  title={Chapter overview},
  fonttitle=\bfseries
}

\newtcolorbox{threedictionaries}{
  colback=DictionaryBlue!4,
  colframe=DictionaryBlue,
  title={The three dictionaries},
  fonttitle=\bfseries
}

\newtcolorbox{surfacelaboratory}{
  colback=SurfaceGreen!5,
  colframe=SurfaceGreen,
  title={Surface laboratory},
  fonttitle=\bfseries
}

\newtcolorbox{warningbox}{
  colback=WarningRed!4,
  colframe=WarningRed,
  title={Warning},
  fonttitle=\bfseries
}

\newtcolorbox{guidingquestion}{
  colback=ExampleGold!5,
  colframe=ExampleGold,
  title={Guiding question},
  fonttitle=\bfseries
}

\newtcolorbox{pauseandthink}{
  colback=PausePurple!5,
  colframe=PausePurple,
  title={Pause and think},
  fonttitle=\bfseries
}

\newtcolorbox{prediction}{
  colback=PredictionTeal!5,
  colframe=PredictionTeal,
  title={Prediction},
  fonttitle=\bfseries
}

\newtcolorbox{dictionarytable}{
  colback=DictionaryBlue!3,
  colframe=DictionaryBlue,
  title={Dictionary summary},
  fonttitle=\bfseries
}

\newtcolorbox{agcheatsheet}{
  colback=CheatsheetSlate!5,
  colframe=CheatsheetSlate,
  title={Algebraic geometry cheatsheet},
  fonttitle=\bfseries
}

% --------------------------------------------------------------------------
% Theorem environments
% --------------------------------------------------------------------------

\theoremstyle{plain}
\newtheorem{theorem}{Theorem}[chapter]
\newtheorem{proposition}[theorem]{Proposition}
\newtheorem{lemma}[theorem]{Lemma}
\newtheorem{corollary}[theorem]{Corollary}

\theoremstyle{definition}
\newtheorem{definition}[theorem]{Definition}
\newtheorem{construction}[theorem]{Construction}
\newtheorem{example}[theorem]{Example}
\newtheorem{exercise}[theorem]{Exercise}
\newtheorem{question}[theorem]{Question}

\theoremstyle{remark}
\newtheorem{remark}[theorem]{Remark}
\newtheorem{observation}[theorem]{Observation}
\newtheorem{convention}[theorem]{Convention}
\newtheorem{notation}[theorem]{Notation}

% Unnumbered variants.
\newtheorem*{theorem*}{Theorem}
\newtheorem*{proposition*}{Proposition}
\newtheorem*{definition*}{Definition}
\newtheorem*{example*}{Example}
\newtheorem*{remark*}{Remark}

% --------------------------------------------------------------------------
% Exercises and solutions
% --------------------------------------------------------------------------

\newcounter{exerciseblock}[chapter]
\renewcommand{\theexerciseblock}{\thechapter.\arabic{exerciseblock}}

\NewDocumentEnvironment{exercises}{O{Exercises}}
{
  \section*{#1}
  \addcontentsline{toc}{section}{#1}
  \begin{enumerate}[label=\textbf{\thechapter.\arabic*.},wide=0pt]
}
{
  \end{enumerate}
}

\newif\ifshowsolutions
\showsolutionsfalse
% To compile instructor solutions, change the previous line to:
% \showsolutionstrue

\NewDocumentEnvironment{solution}{}
{
  \ifshowsolutions
    \par\smallskip
    \noindent\textbf{Solution. }\ignorespaces
  \else
    \expandafter\comment
  \fi
}
{
  \ifshowsolutions
    \par\smallskip
  \else
    \expandafter\endcomment
  \fi
}

% --------------------------------------------------------------------------
% Cross-references
% --------------------------------------------------------------------------

\usepackage{hyperref}

\hypersetup{
  colorlinks=true,
  linkcolor=DictionaryBlue,
  citecolor=SurfaceGreen,
  urlcolor=WarningRed,
  pdfauthor={},
  pdftitle={The Three Dictionaries of Toric Geometry},
  pdfsubject={Notes for Renzo's Fall 2026 course at Colorado State University},
  bookmarksopen=true,
  bookmarksnumbered=true
}

\usepackage[nameinlink,noabbrev,capitalise]{cleveref}

\crefname{theorem}{theorem}{theorems}
\crefname{proposition}{proposition}{propositions}
\crefname{lemma}{lemma}{lemmas}
\crefname{corollary}{corollary}{corollaries}
\crefname{definition}{definition}{definitions}
\crefname{construction}{construction}{constructions}
\crefname{example}{example}{examples}
\crefname{exercise}{exercise}{exercises}
\crefname{question}{question}{questions}
\crefname{remark}{remark}{remarks}
\crefname{observation}{observation}{observations}
\crefname{convention}{convention}{conventions}
\crefname{notation}{notation}{notations}

% --------------------------------------------------------------------------
% Bibliography
% --------------------------------------------------------------------------

\usepackage[
  backend=biber,
  style=alphabetic,
  sorting=nyt,
  maxbibnames=99
]{biblatex}

% The present notes use chapter-level prose references rather than citation
% commands.  This guard keeps the project ready for a future bibliography
% without requiring an empty references.bib file in the current release.
\IfFileExists{references.bib}{\addbibresource{references.bib}}{}

% --------------------------------------------------------------------------
% Index
% --------------------------------------------------------------------------

\usepackage{makeidx}
\makeindex

% --------------------------------------------------------------------------
% Drafting tools
% --------------------------------------------------------------------------

% Margin notes can be used during development:
%     \marginnote{Check this proof.}
\usepackage{marginnote}

\newcommand{\todo}[1]{%
  \marginpar{\raggedright\scriptsize\color{WarningRed}\textbf{TODO:} #1}%
}

% Set to true while drafting to display labels in the margins.
\newif\ifshowlabels
\showlabelsfalse

\ifshowlabels
  \usepackage{showkeys}
\fi

% --------------------------------------------------------------------------
% Miscellaneous formatting
% --------------------------------------------------------------------------

\allowdisplaybreaks
\setcounter{tocdepth}{2}
\setcounter{secnumdepth}{3}

\renewcommand{\qedsymbol}{\(\square\)}

% Prevent isolated lines at page boundaries.
\widowpenalty=10000
\clubpenalty=10000
\displaywidowpenalty=10000

%     % macros.tex
% Mathematical notation for graduate notes on toric varieties.
%
% The conventions are chosen to stay reasonably close to
% Cox--Little--Schenck, Toric Varieties.
%
% Load this file after preamble.tex:
%     % preamble.tex
% Shared package imports and environment definitions for the course notes.
% This file is intended to be loaded from main.tex via
%     \input{preamble}

% --------------------------------------------------------------------------
% Encoding, fonts, and typography
% --------------------------------------------------------------------------

\usepackage[T1]{fontenc}
\usepackage[utf8]{inputenc} % Harmless on modern LaTeX; useful on older systems.
\usepackage{lmodern}
\usepackage{microtype}

% --------------------------------------------------------------------------
% Page layout and document structure
% --------------------------------------------------------------------------

\usepackage[
  letterpaper,
  inner=1.25in,
  outer=1.00in,
  top=1.00in,
  bottom=1.10in,
  headheight=15pt
]{geometry}

\usepackage{setspace}
\setstretch{1.05}

\usepackage{emptypage}
\usepackage{fancyhdr}

\pagestyle{fancy}
\fancyhf{}
\fancyhead[LE]{\small\thepage}
\fancyhead[RO]{\small\thepage}
\fancyhead[LO]{\small\nouppercase{\rightmark}}
\fancyhead[RE]{\small\nouppercase{\leftmark}}
\renewcommand{\headrulewidth}{0.4pt}
\renewcommand{\footrulewidth}{0pt}

\fancypagestyle{plain}{
  \fancyhf{}
  \fancyfoot[C]{\small\thepage}
  \renewcommand{\headrulewidth}{0pt}
}

\usepackage{titlesec}

\titleformat{\chapter}[display]
  {\normalfont\bfseries}
  {\Large Chapter \thechapter}
  {0.5ex}
  {\Huge}

\titleformat{name=\chapter,numberless}[display]
  {\normalfont\bfseries}
  {}
  {0pt}
  {\Huge}

\titlespacing*{\chapter}{0pt}{-10pt}{30pt}

% --------------------------------------------------------------------------
% Mathematics
% --------------------------------------------------------------------------

\usepackage{amsmath}
\usepackage{amssymb}
\usepackage{amsthm}
\usepackage{mathtools}
\usepackage{mathrsfs}
\usepackage{bm}
\usepackage{dsfont}

\numberwithin{equation}{chapter}

% --------------------------------------------------------------------------
% Lists, tables, and general utilities
% --------------------------------------------------------------------------

\usepackage{enumitem}
\setlist{
  itemsep=2pt,
  topsep=4pt,
  parsep=0pt
}

\usepackage{booktabs}
\usepackage{array}
\usepackage{longtable}
\usepackage{tabularx}
\usepackage{multicol}
\usepackage{ragged2e}

\usepackage{xcolor}
\usepackage{xparse}
\usepackage{etoolbox}
\usepackage{ifthen}
\usepackage{comment}

% --------------------------------------------------------------------------
% Graphics, diagrams, and commutative diagrams
% --------------------------------------------------------------------------

\usepackage{graphicx}
\graphicspath{{figures/}}

\usepackage{tikz}
\usetikzlibrary{
  arrows.meta,
  calc,
  decorations.markings,
  intersections,
  matrix,
  patterns,
  positioning,
  shapes.geometric
}

\usepackage{tikz-cd}

% Basic TikZ styles for fans, lattices, and polytopes.
\tikzset{
  lattice point/.style={
    circle,
    fill,
    inner sep=1.4pt
  },
  ray/.style={
    -{Latex[length=2.2mm]},
    thick
  },
  fan cone/.style={
    fill=black!8,
    draw=black!55
  },
  polytope/.style={
    thick
  },
  toric divisor/.style={
    thick,
    dashed
  }
}

% --------------------------------------------------------------------------
% Color palette
% --------------------------------------------------------------------------

\definecolor{DictionaryBlue}{HTML}{1F4E79}
\definecolor{SurfaceGreen}{HTML}{3F6B4F}
\definecolor{WarningRed}{HTML}{8B2E2E}
\definecolor{ExampleGold}{HTML}{8A641B}
\definecolor{LightGray}{HTML}{F4F4F4}
\definecolor{PastelYellow}{HTML}{FFF8CC}
\definecolor{PausePurple}{HTML}{6F5A8A}
\definecolor{PredictionTeal}{HTML}{287D7A}
\definecolor{CheatsheetSlate}{HTML}{536878}

% --------------------------------------------------------------------------
% Boxed pedagogical environments
% --------------------------------------------------------------------------

\usepackage[most]{tcolorbox}

\tcbset{
  enhanced,
  breakable,
  boxrule=0.7pt,
  arc=1.5mm,
  left=2mm,
  right=2mm,
  top=1.5mm,
  bottom=1.5mm,
  before skip=8pt,
  after skip=8pt
}

\newtcolorbox{chapteroverview}{
  colback=PastelYellow,
  colframe=ExampleGold!70!black,
  title={Chapter overview},
  fonttitle=\bfseries
}

\newtcolorbox{threedictionaries}{
  colback=DictionaryBlue!4,
  colframe=DictionaryBlue,
  title={The three dictionaries},
  fonttitle=\bfseries
}

\newtcolorbox{surfacelaboratory}{
  colback=SurfaceGreen!5,
  colframe=SurfaceGreen,
  title={Surface laboratory},
  fonttitle=\bfseries
}

\newtcolorbox{warningbox}{
  colback=WarningRed!4,
  colframe=WarningRed,
  title={Warning},
  fonttitle=\bfseries
}

\newtcolorbox{guidingquestion}{
  colback=ExampleGold!5,
  colframe=ExampleGold,
  title={Guiding question},
  fonttitle=\bfseries
}

\newtcolorbox{pauseandthink}{
  colback=PausePurple!5,
  colframe=PausePurple,
  title={Pause and think},
  fonttitle=\bfseries
}

\newtcolorbox{prediction}{
  colback=PredictionTeal!5,
  colframe=PredictionTeal,
  title={Prediction},
  fonttitle=\bfseries
}

\newtcolorbox{dictionarytable}{
  colback=DictionaryBlue!3,
  colframe=DictionaryBlue,
  title={Dictionary summary},
  fonttitle=\bfseries
}

\newtcolorbox{agcheatsheet}{
  colback=CheatsheetSlate!5,
  colframe=CheatsheetSlate,
  title={Algebraic geometry cheatsheet},
  fonttitle=\bfseries
}

% --------------------------------------------------------------------------
% Theorem environments
% --------------------------------------------------------------------------

\theoremstyle{plain}
\newtheorem{theorem}{Theorem}[chapter]
\newtheorem{proposition}[theorem]{Proposition}
\newtheorem{lemma}[theorem]{Lemma}
\newtheorem{corollary}[theorem]{Corollary}

\theoremstyle{definition}
\newtheorem{definition}[theorem]{Definition}
\newtheorem{construction}[theorem]{Construction}
\newtheorem{example}[theorem]{Example}
\newtheorem{exercise}[theorem]{Exercise}
\newtheorem{question}[theorem]{Question}

\theoremstyle{remark}
\newtheorem{remark}[theorem]{Remark}
\newtheorem{observation}[theorem]{Observation}
\newtheorem{convention}[theorem]{Convention}
\newtheorem{notation}[theorem]{Notation}

% Unnumbered variants.
\newtheorem*{theorem*}{Theorem}
\newtheorem*{proposition*}{Proposition}
\newtheorem*{definition*}{Definition}
\newtheorem*{example*}{Example}
\newtheorem*{remark*}{Remark}

% --------------------------------------------------------------------------
% Exercises and solutions
% --------------------------------------------------------------------------

\newcounter{exerciseblock}[chapter]
\renewcommand{\theexerciseblock}{\thechapter.\arabic{exerciseblock}}

\NewDocumentEnvironment{exercises}{O{Exercises}}
{
  \section*{#1}
  \addcontentsline{toc}{section}{#1}
  \begin{enumerate}[label=\textbf{\thechapter.\arabic*.},wide=0pt]
}
{
  \end{enumerate}
}

\newif\ifshowsolutions
\showsolutionsfalse
% To compile instructor solutions, change the previous line to:
% \showsolutionstrue

\NewDocumentEnvironment{solution}{}
{
  \ifshowsolutions
    \par\smallskip
    \noindent\textbf{Solution. }\ignorespaces
  \else
    \expandafter\comment
  \fi
}
{
  \ifshowsolutions
    \par\smallskip
  \else
    \expandafter\endcomment
  \fi
}

% --------------------------------------------------------------------------
% Cross-references
% --------------------------------------------------------------------------

\usepackage{hyperref}

\hypersetup{
  colorlinks=true,
  linkcolor=DictionaryBlue,
  citecolor=SurfaceGreen,
  urlcolor=WarningRed,
  pdfauthor={},
  pdftitle={The Three Dictionaries of Toric Geometry},
  pdfsubject={Notes for Renzo's Fall 2026 course at Colorado State University},
  bookmarksopen=true,
  bookmarksnumbered=true
}

\usepackage[nameinlink,noabbrev,capitalise]{cleveref}

\crefname{theorem}{theorem}{theorems}
\crefname{proposition}{proposition}{propositions}
\crefname{lemma}{lemma}{lemmas}
\crefname{corollary}{corollary}{corollaries}
\crefname{definition}{definition}{definitions}
\crefname{construction}{construction}{constructions}
\crefname{example}{example}{examples}
\crefname{exercise}{exercise}{exercises}
\crefname{question}{question}{questions}
\crefname{remark}{remark}{remarks}
\crefname{observation}{observation}{observations}
\crefname{convention}{convention}{conventions}
\crefname{notation}{notation}{notations}

% --------------------------------------------------------------------------
% Bibliography
% --------------------------------------------------------------------------

\usepackage[
  backend=biber,
  style=alphabetic,
  sorting=nyt,
  maxbibnames=99
]{biblatex}

% The present notes use chapter-level prose references rather than citation
% commands.  This guard keeps the project ready for a future bibliography
% without requiring an empty references.bib file in the current release.
\IfFileExists{references.bib}{\addbibresource{references.bib}}{}

% --------------------------------------------------------------------------
% Index
% --------------------------------------------------------------------------

\usepackage{makeidx}
\makeindex

% --------------------------------------------------------------------------
% Drafting tools
% --------------------------------------------------------------------------

% Margin notes can be used during development:
%     \marginnote{Check this proof.}
\usepackage{marginnote}

\newcommand{\todo}[1]{%
  \marginpar{\raggedright\scriptsize\color{WarningRed}\textbf{TODO:} #1}%
}

% Set to true while drafting to display labels in the margins.
\newif\ifshowlabels
\showlabelsfalse

\ifshowlabels
  \usepackage{showkeys}
\fi

% --------------------------------------------------------------------------
% Miscellaneous formatting
% --------------------------------------------------------------------------

\allowdisplaybreaks
\setcounter{tocdepth}{2}
\setcounter{secnumdepth}{3}

\renewcommand{\qedsymbol}{\(\square\)}

% Prevent isolated lines at page boundaries.
\widowpenalty=10000
\clubpenalty=10000
\displaywidowpenalty=10000

%     % macros.tex
% Mathematical notation for graduate notes on toric varieties.
%
% The conventions are chosen to stay reasonably close to
% Cox--Little--Schenck, Toric Varieties.
%
% Load this file after preamble.tex:
%     \input{preamble}
%     \input{macros}

% ==========================================================================
% General conventions
% ==========================================================================

% Base field.  The default is an algebraically closed field k.
\providecommand{\Bbbk}{\mathbb{k}}

% Standard number systems.
\newcommand{\NN}{\mathbb{N}}
\newcommand{\ZZ}{\mathbb{Z}}
\newcommand{\QQ}{\mathbb{Q}}
\newcommand{\RR}{\mathbb{R}}
\newcommand{\CC}{\mathbb{C}}

% Affine and projective space.
\newcommand{\AAA}{\mathbb{A}}
\newcommand{\PP}{\mathbb{P}}

% Common script and calligraphic letters.
\newcommand{\cA}{\mathcal{A}}
\newcommand{\cB}{\mathcal{B}}
\newcommand{\cC}{\mathcal{C}}
\newcommand{\cD}{\mathcal{D}}
\newcommand{\cE}{\mathcal{E}}
\newcommand{\cF}{\mathcal{F}}
\newcommand{\cG}{\mathcal{G}}
\newcommand{\cH}{\mathcal{H}}
\newcommand{\cI}{\mathcal{I}}
\newcommand{\cJ}{\mathcal{J}}
\newcommand{\cK}{\mathcal{K}}
\newcommand{\cL}{\mathcal{L}}
\newcommand{\cM}{\mathcal{M}}
\newcommand{\cN}{\mathcal{N}}
\newcommand{\cO}{\mathcal{O}}
\newcommand{\cP}{\mathcal{P}}
\newcommand{\cQ}{\mathcal{Q}}
\newcommand{\cR}{\mathcal{R}}
\newcommand{\cS}{\mathcal{S}}
\newcommand{\cT}{\mathcal{T}}
\newcommand{\cU}{\mathcal{U}}
\newcommand{\cV}{\mathcal{V}}
\newcommand{\cW}{\mathcal{W}}
\newcommand{\cX}{\mathcal{X}}
\newcommand{\cY}{\mathcal{Y}}
\newcommand{\cZ}{\mathcal{Z}}

% Fraktur letters used occasionally for ideals and Lie algebras.
\newcommand{\fm}{\mathfrak{m}}
\newcommand{\fp}{\mathfrak{p}}
\newcommand{\fq}{\mathfrak{q}}
\newcommand{\fg}{\mathfrak{g}}
\newcommand{\ft}{\mathfrak{t}}

% Delimiters.
\DeclarePairedDelimiter{\paren}{(}{)}
\DeclarePairedDelimiter{\bracket}{[}{]}
\DeclarePairedDelimiter{\set}{\{}{\}}
\DeclarePairedDelimiter{\abs}{\lvert}{\rvert}
\DeclarePairedDelimiter{\norm}{\lVert}{\rVert}
\DeclarePairedDelimiter{\angles}{\langle}{\rangle}
\DeclarePairedDelimiter{\floor}{\lfloor}{\rfloor}
\DeclarePairedDelimiter{\ceil}{\lceil}{\rceil}

% Divisibility and restrictions.
\newcommand{\divides}{\mathrel{\mid}}
\newcommand{\ndivides}{\mathrel{\nmid}}
\newcommand{\restr}[2]{\left.#1\right|_{#2}}

% ==========================================================================
% Operators
% ==========================================================================

\DeclareMathOperator{\Spec}{Spec}
\DeclareMathOperator{\Proj}{Proj}
\DeclareMathOperator{\MaxSpec}{MaxSpec}
\DeclareMathOperator{\Hom}{Hom}
\DeclareMathOperator{\End}{End}
\DeclareMathOperator{\Aut}{Aut}
\DeclareMathOperator{\Isom}{Isom}
\DeclareMathOperator{\Mor}{Mor}
\DeclareMathOperator{\Ext}{Ext}
\DeclareMathOperator{\Tor}{Tor}
\DeclareMathOperator{\Pic}{Pic}
\DeclareMathOperator{\Cl}{Cl}
\DeclareMathOperator{\Div}{Div}
\DeclareMathOperator{\CaDiv}{CaDiv}
\DeclareMathOperator{\WDiv}{WDiv}
\DeclareMathOperator{\Prin}{Prin}
\DeclareMathOperator{\Supp}{Supp}
\DeclareMathOperator{\ord}{ord}
\DeclareMathOperator{\codim}{codim}
\DeclareMathOperator{\rank}{rank}
\DeclareMathOperator{\Span}{span}
\DeclareMathOperator{\cone}{Cone}
\DeclareMathOperator{\Conv}{Conv}
\DeclareMathOperator{\relint}{Relint}
\DeclareMathOperator{\Aff}{Aff}
\DeclareMathOperator{\vol}{Vol}
\DeclareMathOperator{\mult}{mult}
\DeclareMathOperator{\length}{length}
\DeclareMathOperator{\im}{im}
\DeclareMathOperator{\coker}{coker}
\DeclareMathOperator{\kerop}{ker}
\DeclareMathOperator{\id}{id}
\DeclareMathOperator{\diag}{diag}
\DeclareMathOperator{\Gal}{Gal}
\DeclareMathOperator{\Frac}{Frac}
\DeclareMathOperator{\trdeg}{trdeg}
\DeclareMathOperator{\Sing}{Sing}
\DeclareMathOperator{\Reg}{Reg}
\DeclareMathOperator{\Exc}{Exc}
\DeclareMathOperator{\Bl}{Bl}
\DeclareMathOperator{\NE}{NE}
\DeclareMathOperator{\Nef}{Nef}
\DeclareMathOperator{\Eff}{Eff}
\DeclareMathOperator{\Mov}{Mov}
\DeclareMathOperator{\Amp}{Amp}
\DeclareMathOperator{\Star}{Star}
\DeclareMathOperator{\Link}{Link}
\DeclareMathOperator{\Homog}{Homog}
\DeclareMathOperator{\Ann}{Ann}
\DeclareMathOperator{\ch}{ch}
\DeclareMathOperator{\td}{td}

% Prefer \ker for ordinary kernels; \kerop is supplied only when an operator
% with explicit limits is useful.
\renewcommand{\ker}{\operatorname{ker}}

% ==========================================================================
% Lattices, tori, characters, and one-parameter subgroups
% ==========================================================================

% CLS convention:
%   N is the lattice of one-parameter subgroups,
%   M = Hom_Z(N,Z) is the character lattice.
\newcommand{\Nlattice}{N}
\newcommand{\Mlattice}{M}

\newcommand{\NR}{N_{\RR}}
\newcommand{\NQ}{N_{\QQ}}
\newcommand{\MR}{M_{\RR}}
\newcommand{\MQ}{M_{\QQ}}

\newcommand{\NRbar}{\overline{N}_{\RR}}
\newcommand{\MRbar}{\overline{M}_{\RR}}

% Tensor extension of an arbitrary lattice.
\newcommand{\LR}[1]{#1_{\RR}}
\newcommand{\LQ}[1]{#1_{\QQ}}

% Dual lattice and dual vector space.
\newcommand{\dual}[1]{#1^{\vee}}
\newcommand{\ZZdual}[1]{\Hom_{\ZZ}(#1,\ZZ)}
\newcommand{\RRdual}[1]{\Hom_{\RR}(#1,\RR)}

% Pairing between M and N.
\newcommand{\pair}[2]{\left\langle #1,#2\right\rangle}

% Algebraic torus associated to N.
\newcommand{\TN}{T_N}
\newcommand{\TM}{T_M}
\newcommand{\torus}{T_N}
\newcommand{\Gm}{\mathbb{G}_{m}}
\newcommand{\Ga}{\mathbb{G}_{a}}

% Characters and one-parameter subgroups.
\newcommand{\character}[1]{\chi^{#1}}
\newcommand{\oneps}[1]{\lambda^{#1}}
\newcommand{\chiM}{\chi^m}
\newcommand{\lambdaU}{\lambda^u}

% Character and cocharacter groups.
\newcommand{\Xstar}[1]{X^{*}(#1)}
\newcommand{\Xsubstar}[1]{X_{*}(#1)}

% Group algebra / Laurent polynomial algebra.
\newcommand{\groupalg}[1]{\Bbbk[#1]}
\newcommand{\laurentring}[1]{\Bbbk[#1^{\pm 1}]}

% Weight spaces and gradings.
\newcommand{\weightspace}[2]{#1_{#2}}
\newcommand{\gradedpiece}[2]{#1_{#2}}
\newcommand{\wt}{\operatorname{wt}}

% ==========================================================================
% Semigroups, cones, and affine toric varieties
% ==========================================================================

% Semigroup generated by a subset A.
\newcommand{\semigroup}[1]{\NN #1}
\newcommand{\ZA}{\ZZ A}
\newcommand{\NA}{\NN A}

% Semigroup algebra.
\newcommand{\semigroupalg}[1]{\Bbbk[#1]}

% Dual cone and orthogonal complement.
\newcommand{\dualcone}[1]{#1^{\vee}}
\newcommand{\perpcone}[1]{#1^{\perp}}

% Cone attached to generators.
\newcommand{\Cone}[1]{\operatorname{Cone}(#1)}
\newcommand{\pos}[1]{\operatorname{pos}(#1)}

% Standard affine toric variety and semigroup associated to a cone.
\newcommand{\Ssigma}{S_{\sigma}}
\newcommand{\Usigma}{U_{\sigma}}
\newcommand{\Uface}[1]{U_{#1}}
\newcommand{\Xsigma}{X_{\sigma}}

% Generic versions.
\newcommand{\Scone}[1]{S_{#1}}
\newcommand{\Ucone}[1]{U_{#1}}
\newcommand{\Xcone}[1]{X_{#1}}

% Distinguished monomial.
\newcommand{\xmon}[1]{x^{#1}}

% Toric ideal associated to A.
\newcommand{\IA}{I_A}
\newcommand{\Iconfig}[1]{I_{#1}}

% Saturation.
\newcommand{\sat}[1]{\overline{#1}}
\newcommand{\saturation}[2]{#1:#2^{\infty}}

% Faces.
\newcommand{\faceleq}{\preceq}
\newcommand{\facet}{\text{facet}}

% ==========================================================================
% Fans and global toric varieties
% ==========================================================================

\newcommand{\fan}{\Sigma}
\newcommand{\fanp}{\Sigma'}
\newcommand{\fanq}{\Sigma''}

\newcommand{\Xfan}{X_{\Sigma}}
\newcommand{\Xfanp}{X_{\Sigma'}}
\newcommand{\Xof}[1]{X_{#1}}

% Cones of a fixed dimension.
\newcommand{\fancones}[1]{\Sigma(#1)}
\newcommand{\fanrays}{\Sigma(1)}
\newcommand{\fanmax}{\Sigma_{\max}}

% Support of a fan.
\newcommand{\fansupp}[1]{\lvert #1\rvert}

% Ray generator.  For rho in Sigma(1), u_rho denotes the primitive generator.
\newcommand{\urho}{u_{\rho}}
\newcommand{\uray}[1]{u_{#1}}
\newcommand{\vrho}{v_{\rho}}
\newcommand{\vray}[1]{v_{#1}}

% Orbit and orbit closure associated to a cone.
\newcommand{\Osigma}{O(\sigma)}
\newcommand{\Vsigma}{V(\sigma)}
\newcommand{\Ocone}[1]{O(#1)}
\newcommand{\Vcone}[1]{V(#1)}

% Star of a cone and quotient fan.
\newcommand{\starsigma}{\operatorname{Star}(\sigma)}
\newcommand{\starcone}[1]{\operatorname{Star}(#1)}

% Refinements and subdivisions.
\newcommand{\refines}{\succ}
\newcommand{\subdivision}[2]{#1 \to #2}

% ==========================================================================
% Polytopes and projective toric varieties
% ==========================================================================

\newcommand{\polytope}{P}
\newcommand{\polytopeQ}{Q}
\newcommand{\DeltaPoly}{\Delta}

\newcommand{\conv}[1]{\operatorname{Conv}(#1)}
\newcommand{\normalfan}[1]{\Sigma_{#1}}
\newcommand{\face}[2]{F_{#1}(#2)}

% Lattice points of a polytope.
\newcommand{\latticepoints}[1]{#1\cap M}
\newcommand{\numlatticepoints}[1]{\#(#1\cap M)}

% Polyhedron associated to a torus-invariant divisor.
\newcommand{\PD}{P_D}
\newcommand{\Pdiv}[1]{P_{#1}}

% Projective toric variety associated to a configuration or polytope.
\newcommand{\XA}{X_A}
\newcommand{\XP}{X_P}
\newcommand{\Xpoly}[1]{X_{#1}}

% Normal fan and support function.
\newcommand{\psiD}{\psi_D}
\newcommand{\supportfunction}[1]{\psi_{#1}}

% ==========================================================================
% Divisors, line bundles, and sections
% ==========================================================================

% Torus-invariant prime divisor associated to rho.
\newcommand{\Drho}{D_{\rho}}
\newcommand{\Dray}[1]{D_{#1}}

% Generic torus-invariant divisor.
\newcommand{\Dsum}{\sum_{\rho\in\Sigma(1)} a_{\rho}D_{\rho}}
\newcommand{\Dtoric}{D=\sum_{\rho\in\Sigma(1)}a_{\rho}D_{\rho}}

% Principal divisor of a character.
\newcommand{\divchi}[1]{\operatorname{div}(\chi^{#1})}

% Divisor and line-bundle notation.
\newcommand{\OX}{\mathcal{O}_X}
\newcommand{\OD}{\mathcal{O}_X(D)}
\newcommand{\Odiv}[1]{\mathcal{O}_X(#1)}
\newcommand{\Ldiv}[1]{\mathcal{L}(#1)}

% Global sections and cohomology.
\newcommand{\Hzero}[2]{H^{0}\!\left(#1,#2\right)}
\newcommand{\Hi}[3]{H^{#1}\!\left(#2,#3\right)}
\newcommand{\hzero}[2]{h^{0}\!\left(#1,#2\right)}
\newcommand{\hi}[3]{h^{#1}\!\left(#2,#3\right)}

% Canonical and anticanonical divisors.
\newcommand{\KX}{K_X}
\newcommand{\minusKX}{-K_X}
\newcommand{\omegaX}{\omega_X}

% Boundary divisor.
\newcommand{\boundaryX}{\partial X}
\newcommand{\toricboundary}{\sum_{\rho\in\Sigma(1)}D_{\rho}}

% Divisor class groups and numerical spaces.
\newcommand{\ClX}{\Cl(X)}
\newcommand{\PicX}{\Pic(X)}
\newcommand{\NSX}{\operatorname{NS}(X)}
\newcommand{\NOne}[1]{N^{1}(#1)_{\RR}}
\newcommand{\Ncurve}[1]{N_{1}(#1)_{\RR}}

% Section ring.
\newcommand{\sectionring}[2]{R(#1,#2)}
\newcommand{\RofD}{R(X,D)}

% ==========================================================================
% Cox rings and homogeneous coordinates
% ==========================================================================

\newcommand{\Cox}[1]{\operatorname{Cox}(#1)}
\newcommand{\CoxX}{\operatorname{Cox}(X_{\Sigma})}
\newcommand{\Scox}{S}
\newcommand{\Bfan}{B(\Sigma)}
\newcommand{\Zfan}{Z(\Sigma)}
\newcommand{\Uhat}{\widehat{X}}
\newcommand{\Gfan}{G_{\Sigma}}

% Cox variables.
\newcommand{\xrho}{x_{\rho}}
\newcommand{\xhat}[1]{\widehat{x}_{#1}}

% Homogeneous piece of degree alpha.
\newcommand{\Salpha}{S_{\alpha}}
\newcommand{\Sdegree}[1]{S_{#1}}

% ==========================================================================
% Toric surfaces and blowups
% ==========================================================================

% Hirzebruch surfaces.
\newcommand{\Hirz}[1]{\mathbb{F}_{#1}}
\newcommand{\Fn}{\mathbb{F}_{n}}

% Blowups.
\newcommand{\Blp}[2]{\operatorname{Bl}_{#2}(#1)}
\newcommand{\blowupmap}{\pi\colon \widetilde{X}\to X}
\newcommand{\Xtilde}{\widetilde{X}}
\newcommand{\Dtilde}{\widetilde{D}}
\newcommand{\Eexc}{E}

% Standard divisor classes on blowups of P^2.
\newcommand{\Hclass}{H}
\newcommand{\Eclass}[1]{E_{#1}}

% Self-intersection.
\newcommand{\selfint}[1]{#1^{2}}
\newcommand{\intersection}[2]{#1\cdot #2}

% Cyclically ordered primitive ray generators on a surface.
\newcommand{\vi}{v_i}
\newcommand{\vip}{v_{i+1}}
\newcommand{\vim}{v_{i-1}}

% ==========================================================================
% Singularities, resolutions, and birational geometry
% ==========================================================================

\newcommand{\resolution}{\pi\colon \widetilde{X}\to X}
\newcommand{\exceptional}{\operatorname{Exc}(\pi)}

\newcommand{\QCartier}{\QQ\text{-Cartier}}
\newcommand{\Qfactorial}{\QQ\text{-factorial}}
\newcommand{\QGorenstein}{\QQ\text{-Gorenstein}}

\newcommand{\MMP}{\mathrm{MMP}}
\newcommand{\NEbar}[1]{\overline{\operatorname{NE}}(#1)}
\newcommand{\Effbar}[1]{\overline{\operatorname{Eff}}(#1)}
\newcommand{\Movbar}[1]{\overline{\operatorname{Mov}}(#1)}

% Discrepancy notation.
\newcommand{\discrepancy}[2]{a(#1,#2)}

% ==========================================================================
% Intersection theory and Chow groups
% ==========================================================================

\newcommand{\Astar}[1]{A^{*}(#1)}
\newcommand{\Aunder}[1]{A_{*}(#1)}
\newcommand{\Ak}[2]{A^{#1}(#2)}
\newcommand{\Acycle}[2]{A_{#1}(#2)}

\newcommand{\chowring}[1]{A^{*}(#1)}
\newcommand{\cycleclass}[1]{[#1]}
\newcommand{\degree}{\operatorname{deg}}

% Stanley--Reisner and linear-relation ideals.
\newcommand{\ISR}{I_{\mathrm{SR}}}
\newcommand{\Ilin}{I_{\mathrm{lin}}}

% ==========================================================================
% Topology and compact torus notation
% ==========================================================================

\newcommand{\compacttorus}{(S^{1})^{n}}
\newcommand{\momentmap}{\mu}
\newcommand{\momentpolytope}{\Delta}

\newcommand{\piOne}[1]{\pi_{1}(#1)}
\newcommand{\Euler}[1]{\chi(#1)}

% ==========================================================================
% Morphisms and categorical notation
% ==========================================================================

\newcommand{\longto}{\longrightarrow}
\newcommand{\longfrom}{\longleftarrow}
\newcommand{\into}{\hookrightarrow}
\newcommand{\onto}{\twoheadrightarrow}
\newcommand{\isoto}{\xrightarrow{\sim}}
\newcommand{\birat}{\dashrightarrow}

\newcommand{\map}[3]{#1\colon #2\longrightarrow #3}
\newcommand{\rmap}[3]{#1\colon #2\dashrightarrow #3}

% Pullback and pushforward.
\newcommand{\pullback}[1]{#1^{*}}
\newcommand{\pushforward}[1]{#1_{*}}

% ==========================================================================
% Frequently used equations and constructions
% ==========================================================================

% The standard exact sequence for torus-invariant divisors:
%     M -> Z^{Sigma(1)} -> Cl(X_Sigma) -> 0.
\newcommand{\divisorsequence}{%
  M \longrightarrow \ZZ^{\Sigma(1)}
    \longrightarrow \Cl(X_{\Sigma}) \longrightarrow 0%
}

% Smooth toric divisor sequence:
\newcommand{\picardsequence}{%
  M \longrightarrow \ZZ^{\Sigma(1)}
    \longrightarrow \Pic(X_{\Sigma}) \longrightarrow 0%
}

% Character divisor formula.
\newcommand{\characterdivisorformula}{%
  \operatorname{div}(\chi^{m})
  =
  \sum_{\rho\in\Sigma(1)}
  \langle m,u_{\rho}\rangle D_{\rho}%
}

% Canonical divisor formula.
\newcommand{\canonicaldivisorformula}{%
  K_{X_{\Sigma}}
  =
  -\sum_{\rho\in\Sigma(1)}D_{\rho}%
}

% ==========================================================================
% Textual abbreviations
% ==========================================================================

\newcommand{\ie}{i.e.\ }
\newcommand{\eg}{e.g.\ }
\newcommand{\cf}{cf.\ }
\newcommand{\resp}{respectively}
\newcommand{\withoutloss}{without loss of generality}
\newcommand{\Withoutloss}{Without loss of generality}
\newcommand{\WLOG}{\text{without loss of generality}}

% ==========================================================================
% Optional notation aliases
% ==========================================================================
%
% These aliases are convenient but somewhat terse.  Keeping them grouped here
% makes it easy to remove them if a notation conflict appears later.

\newcommand{\Sig}{\Sigma}
\newcommand{\sig}{\sigma}
\newcommand{\taut}{\tau}
\newcommand{\rhoRay}{\rho}

% End of macros.tex


% ==========================================================================
% General conventions
% ==========================================================================

% Base field.  The default is an algebraically closed field k.
\providecommand{\Bbbk}{\mathbb{k}}

% Standard number systems.
\newcommand{\NN}{\mathbb{N}}
\newcommand{\ZZ}{\mathbb{Z}}
\newcommand{\QQ}{\mathbb{Q}}
\newcommand{\RR}{\mathbb{R}}
\newcommand{\CC}{\mathbb{C}}

% Affine and projective space.
\newcommand{\AAA}{\mathbb{A}}
\newcommand{\PP}{\mathbb{P}}

% Common script and calligraphic letters.
\newcommand{\cA}{\mathcal{A}}
\newcommand{\cB}{\mathcal{B}}
\newcommand{\cC}{\mathcal{C}}
\newcommand{\cD}{\mathcal{D}}
\newcommand{\cE}{\mathcal{E}}
\newcommand{\cF}{\mathcal{F}}
\newcommand{\cG}{\mathcal{G}}
\newcommand{\cH}{\mathcal{H}}
\newcommand{\cI}{\mathcal{I}}
\newcommand{\cJ}{\mathcal{J}}
\newcommand{\cK}{\mathcal{K}}
\newcommand{\cL}{\mathcal{L}}
\newcommand{\cM}{\mathcal{M}}
\newcommand{\cN}{\mathcal{N}}
\newcommand{\cO}{\mathcal{O}}
\newcommand{\cP}{\mathcal{P}}
\newcommand{\cQ}{\mathcal{Q}}
\newcommand{\cR}{\mathcal{R}}
\newcommand{\cS}{\mathcal{S}}
\newcommand{\cT}{\mathcal{T}}
\newcommand{\cU}{\mathcal{U}}
\newcommand{\cV}{\mathcal{V}}
\newcommand{\cW}{\mathcal{W}}
\newcommand{\cX}{\mathcal{X}}
\newcommand{\cY}{\mathcal{Y}}
\newcommand{\cZ}{\mathcal{Z}}

% Fraktur letters used occasionally for ideals and Lie algebras.
\newcommand{\fm}{\mathfrak{m}}
\newcommand{\fp}{\mathfrak{p}}
\newcommand{\fq}{\mathfrak{q}}
\newcommand{\fg}{\mathfrak{g}}
\newcommand{\ft}{\mathfrak{t}}

% Delimiters.
\DeclarePairedDelimiter{\paren}{(}{)}
\DeclarePairedDelimiter{\bracket}{[}{]}
\DeclarePairedDelimiter{\set}{\{}{\}}
\DeclarePairedDelimiter{\abs}{\lvert}{\rvert}
\DeclarePairedDelimiter{\norm}{\lVert}{\rVert}
\DeclarePairedDelimiter{\angles}{\langle}{\rangle}
\DeclarePairedDelimiter{\floor}{\lfloor}{\rfloor}
\DeclarePairedDelimiter{\ceil}{\lceil}{\rceil}

% Divisibility and restrictions.
\newcommand{\divides}{\mathrel{\mid}}
\newcommand{\ndivides}{\mathrel{\nmid}}
\newcommand{\restr}[2]{\left.#1\right|_{#2}}

% ==========================================================================
% Operators
% ==========================================================================

\DeclareMathOperator{\Spec}{Spec}
\DeclareMathOperator{\Proj}{Proj}
\DeclareMathOperator{\MaxSpec}{MaxSpec}
\DeclareMathOperator{\Hom}{Hom}
\DeclareMathOperator{\End}{End}
\DeclareMathOperator{\Aut}{Aut}
\DeclareMathOperator{\Isom}{Isom}
\DeclareMathOperator{\Mor}{Mor}
\DeclareMathOperator{\Ext}{Ext}
\DeclareMathOperator{\Tor}{Tor}
\DeclareMathOperator{\Pic}{Pic}
\DeclareMathOperator{\Cl}{Cl}
\DeclareMathOperator{\Div}{Div}
\DeclareMathOperator{\CaDiv}{CaDiv}
\DeclareMathOperator{\WDiv}{WDiv}
\DeclareMathOperator{\Prin}{Prin}
\DeclareMathOperator{\Supp}{Supp}
\DeclareMathOperator{\ord}{ord}
\DeclareMathOperator{\codim}{codim}
\DeclareMathOperator{\rank}{rank}
\DeclareMathOperator{\Span}{span}
\DeclareMathOperator{\cone}{Cone}
\DeclareMathOperator{\Conv}{Conv}
\DeclareMathOperator{\relint}{Relint}
\DeclareMathOperator{\Aff}{Aff}
\DeclareMathOperator{\vol}{Vol}
\DeclareMathOperator{\mult}{mult}
\DeclareMathOperator{\length}{length}
\DeclareMathOperator{\im}{im}
\DeclareMathOperator{\coker}{coker}
\DeclareMathOperator{\kerop}{ker}
\DeclareMathOperator{\id}{id}
\DeclareMathOperator{\diag}{diag}
\DeclareMathOperator{\Gal}{Gal}
\DeclareMathOperator{\Frac}{Frac}
\DeclareMathOperator{\trdeg}{trdeg}
\DeclareMathOperator{\Sing}{Sing}
\DeclareMathOperator{\Reg}{Reg}
\DeclareMathOperator{\Exc}{Exc}
\DeclareMathOperator{\Bl}{Bl}
\DeclareMathOperator{\NE}{NE}
\DeclareMathOperator{\Nef}{Nef}
\DeclareMathOperator{\Eff}{Eff}
\DeclareMathOperator{\Mov}{Mov}
\DeclareMathOperator{\Amp}{Amp}
\DeclareMathOperator{\Star}{Star}
\DeclareMathOperator{\Link}{Link}
\DeclareMathOperator{\Homog}{Homog}
\DeclareMathOperator{\Ann}{Ann}
\DeclareMathOperator{\ch}{ch}
\DeclareMathOperator{\td}{td}

% Prefer \ker for ordinary kernels; \kerop is supplied only when an operator
% with explicit limits is useful.
\renewcommand{\ker}{\operatorname{ker}}

% ==========================================================================
% Lattices, tori, characters, and one-parameter subgroups
% ==========================================================================

% CLS convention:
%   N is the lattice of one-parameter subgroups,
%   M = Hom_Z(N,Z) is the character lattice.
\newcommand{\Nlattice}{N}
\newcommand{\Mlattice}{M}

\newcommand{\NR}{N_{\RR}}
\newcommand{\NQ}{N_{\QQ}}
\newcommand{\MR}{M_{\RR}}
\newcommand{\MQ}{M_{\QQ}}

\newcommand{\NRbar}{\overline{N}_{\RR}}
\newcommand{\MRbar}{\overline{M}_{\RR}}

% Tensor extension of an arbitrary lattice.
\newcommand{\LR}[1]{#1_{\RR}}
\newcommand{\LQ}[1]{#1_{\QQ}}

% Dual lattice and dual vector space.
\newcommand{\dual}[1]{#1^{\vee}}
\newcommand{\ZZdual}[1]{\Hom_{\ZZ}(#1,\ZZ)}
\newcommand{\RRdual}[1]{\Hom_{\RR}(#1,\RR)}

% Pairing between M and N.
\newcommand{\pair}[2]{\left\langle #1,#2\right\rangle}

% Algebraic torus associated to N.
\newcommand{\TN}{T_N}
\newcommand{\TM}{T_M}
\newcommand{\torus}{T_N}
\newcommand{\Gm}{\mathbb{G}_{m}}
\newcommand{\Ga}{\mathbb{G}_{a}}

% Characters and one-parameter subgroups.
\newcommand{\character}[1]{\chi^{#1}}
\newcommand{\oneps}[1]{\lambda^{#1}}
\newcommand{\chiM}{\chi^m}
\newcommand{\lambdaU}{\lambda^u}

% Character and cocharacter groups.
\newcommand{\Xstar}[1]{X^{*}(#1)}
\newcommand{\Xsubstar}[1]{X_{*}(#1)}

% Group algebra / Laurent polynomial algebra.
\newcommand{\groupalg}[1]{\Bbbk[#1]}
\newcommand{\laurentring}[1]{\Bbbk[#1^{\pm 1}]}

% Weight spaces and gradings.
\newcommand{\weightspace}[2]{#1_{#2}}
\newcommand{\gradedpiece}[2]{#1_{#2}}
\newcommand{\wt}{\operatorname{wt}}

% ==========================================================================
% Semigroups, cones, and affine toric varieties
% ==========================================================================

% Semigroup generated by a subset A.
\newcommand{\semigroup}[1]{\NN #1}
\newcommand{\ZA}{\ZZ A}
\newcommand{\NA}{\NN A}

% Semigroup algebra.
\newcommand{\semigroupalg}[1]{\Bbbk[#1]}

% Dual cone and orthogonal complement.
\newcommand{\dualcone}[1]{#1^{\vee}}
\newcommand{\perpcone}[1]{#1^{\perp}}

% Cone attached to generators.
\newcommand{\Cone}[1]{\operatorname{Cone}(#1)}
\newcommand{\pos}[1]{\operatorname{pos}(#1)}

% Standard affine toric variety and semigroup associated to a cone.
\newcommand{\Ssigma}{S_{\sigma}}
\newcommand{\Usigma}{U_{\sigma}}
\newcommand{\Uface}[1]{U_{#1}}
\newcommand{\Xsigma}{X_{\sigma}}

% Generic versions.
\newcommand{\Scone}[1]{S_{#1}}
\newcommand{\Ucone}[1]{U_{#1}}
\newcommand{\Xcone}[1]{X_{#1}}

% Distinguished monomial.
\newcommand{\xmon}[1]{x^{#1}}

% Toric ideal associated to A.
\newcommand{\IA}{I_A}
\newcommand{\Iconfig}[1]{I_{#1}}

% Saturation.
\newcommand{\sat}[1]{\overline{#1}}
\newcommand{\saturation}[2]{#1:#2^{\infty}}

% Faces.
\newcommand{\faceleq}{\preceq}
\newcommand{\facet}{\text{facet}}

% ==========================================================================
% Fans and global toric varieties
% ==========================================================================

\newcommand{\fan}{\Sigma}
\newcommand{\fanp}{\Sigma'}
\newcommand{\fanq}{\Sigma''}

\newcommand{\Xfan}{X_{\Sigma}}
\newcommand{\Xfanp}{X_{\Sigma'}}
\newcommand{\Xof}[1]{X_{#1}}

% Cones of a fixed dimension.
\newcommand{\fancones}[1]{\Sigma(#1)}
\newcommand{\fanrays}{\Sigma(1)}
\newcommand{\fanmax}{\Sigma_{\max}}

% Support of a fan.
\newcommand{\fansupp}[1]{\lvert #1\rvert}

% Ray generator.  For rho in Sigma(1), u_rho denotes the primitive generator.
\newcommand{\urho}{u_{\rho}}
\newcommand{\uray}[1]{u_{#1}}
\newcommand{\vrho}{v_{\rho}}
\newcommand{\vray}[1]{v_{#1}}

% Orbit and orbit closure associated to a cone.
\newcommand{\Osigma}{O(\sigma)}
\newcommand{\Vsigma}{V(\sigma)}
\newcommand{\Ocone}[1]{O(#1)}
\newcommand{\Vcone}[1]{V(#1)}

% Star of a cone and quotient fan.
\newcommand{\starsigma}{\operatorname{Star}(\sigma)}
\newcommand{\starcone}[1]{\operatorname{Star}(#1)}

% Refinements and subdivisions.
\newcommand{\refines}{\succ}
\newcommand{\subdivision}[2]{#1 \to #2}

% ==========================================================================
% Polytopes and projective toric varieties
% ==========================================================================

\newcommand{\polytope}{P}
\newcommand{\polytopeQ}{Q}
\newcommand{\DeltaPoly}{\Delta}

\newcommand{\conv}[1]{\operatorname{Conv}(#1)}
\newcommand{\normalfan}[1]{\Sigma_{#1}}
\newcommand{\face}[2]{F_{#1}(#2)}

% Lattice points of a polytope.
\newcommand{\latticepoints}[1]{#1\cap M}
\newcommand{\numlatticepoints}[1]{\#(#1\cap M)}

% Polyhedron associated to a torus-invariant divisor.
\newcommand{\PD}{P_D}
\newcommand{\Pdiv}[1]{P_{#1}}

% Projective toric variety associated to a configuration or polytope.
\newcommand{\XA}{X_A}
\newcommand{\XP}{X_P}
\newcommand{\Xpoly}[1]{X_{#1}}

% Normal fan and support function.
\newcommand{\psiD}{\psi_D}
\newcommand{\supportfunction}[1]{\psi_{#1}}

% ==========================================================================
% Divisors, line bundles, and sections
% ==========================================================================

% Torus-invariant prime divisor associated to rho.
\newcommand{\Drho}{D_{\rho}}
\newcommand{\Dray}[1]{D_{#1}}

% Generic torus-invariant divisor.
\newcommand{\Dsum}{\sum_{\rho\in\Sigma(1)} a_{\rho}D_{\rho}}
\newcommand{\Dtoric}{D=\sum_{\rho\in\Sigma(1)}a_{\rho}D_{\rho}}

% Principal divisor of a character.
\newcommand{\divchi}[1]{\operatorname{div}(\chi^{#1})}

% Divisor and line-bundle notation.
\newcommand{\OX}{\mathcal{O}_X}
\newcommand{\OD}{\mathcal{O}_X(D)}
\newcommand{\Odiv}[1]{\mathcal{O}_X(#1)}
\newcommand{\Ldiv}[1]{\mathcal{L}(#1)}

% Global sections and cohomology.
\newcommand{\Hzero}[2]{H^{0}\!\left(#1,#2\right)}
\newcommand{\Hi}[3]{H^{#1}\!\left(#2,#3\right)}
\newcommand{\hzero}[2]{h^{0}\!\left(#1,#2\right)}
\newcommand{\hi}[3]{h^{#1}\!\left(#2,#3\right)}

% Canonical and anticanonical divisors.
\newcommand{\KX}{K_X}
\newcommand{\minusKX}{-K_X}
\newcommand{\omegaX}{\omega_X}

% Boundary divisor.
\newcommand{\boundaryX}{\partial X}
\newcommand{\toricboundary}{\sum_{\rho\in\Sigma(1)}D_{\rho}}

% Divisor class groups and numerical spaces.
\newcommand{\ClX}{\Cl(X)}
\newcommand{\PicX}{\Pic(X)}
\newcommand{\NSX}{\operatorname{NS}(X)}
\newcommand{\NOne}[1]{N^{1}(#1)_{\RR}}
\newcommand{\Ncurve}[1]{N_{1}(#1)_{\RR}}

% Section ring.
\newcommand{\sectionring}[2]{R(#1,#2)}
\newcommand{\RofD}{R(X,D)}

% ==========================================================================
% Cox rings and homogeneous coordinates
% ==========================================================================

\newcommand{\Cox}[1]{\operatorname{Cox}(#1)}
\newcommand{\CoxX}{\operatorname{Cox}(X_{\Sigma})}
\newcommand{\Scox}{S}
\newcommand{\Bfan}{B(\Sigma)}
\newcommand{\Zfan}{Z(\Sigma)}
\newcommand{\Uhat}{\widehat{X}}
\newcommand{\Gfan}{G_{\Sigma}}

% Cox variables.
\newcommand{\xrho}{x_{\rho}}
\newcommand{\xhat}[1]{\widehat{x}_{#1}}

% Homogeneous piece of degree alpha.
\newcommand{\Salpha}{S_{\alpha}}
\newcommand{\Sdegree}[1]{S_{#1}}

% ==========================================================================
% Toric surfaces and blowups
% ==========================================================================

% Hirzebruch surfaces.
\newcommand{\Hirz}[1]{\mathbb{F}_{#1}}
\newcommand{\Fn}{\mathbb{F}_{n}}

% Blowups.
\newcommand{\Blp}[2]{\operatorname{Bl}_{#2}(#1)}
\newcommand{\blowupmap}{\pi\colon \widetilde{X}\to X}
\newcommand{\Xtilde}{\widetilde{X}}
\newcommand{\Dtilde}{\widetilde{D}}
\newcommand{\Eexc}{E}

% Standard divisor classes on blowups of P^2.
\newcommand{\Hclass}{H}
\newcommand{\Eclass}[1]{E_{#1}}

% Self-intersection.
\newcommand{\selfint}[1]{#1^{2}}
\newcommand{\intersection}[2]{#1\cdot #2}

% Cyclically ordered primitive ray generators on a surface.
\newcommand{\vi}{v_i}
\newcommand{\vip}{v_{i+1}}
\newcommand{\vim}{v_{i-1}}

% ==========================================================================
% Singularities, resolutions, and birational geometry
% ==========================================================================

\newcommand{\resolution}{\pi\colon \widetilde{X}\to X}
\newcommand{\exceptional}{\operatorname{Exc}(\pi)}

\newcommand{\QCartier}{\QQ\text{-Cartier}}
\newcommand{\Qfactorial}{\QQ\text{-factorial}}
\newcommand{\QGorenstein}{\QQ\text{-Gorenstein}}

\newcommand{\MMP}{\mathrm{MMP}}
\newcommand{\NEbar}[1]{\overline{\operatorname{NE}}(#1)}
\newcommand{\Effbar}[1]{\overline{\operatorname{Eff}}(#1)}
\newcommand{\Movbar}[1]{\overline{\operatorname{Mov}}(#1)}

% Discrepancy notation.
\newcommand{\discrepancy}[2]{a(#1,#2)}

% ==========================================================================
% Intersection theory and Chow groups
% ==========================================================================

\newcommand{\Astar}[1]{A^{*}(#1)}
\newcommand{\Aunder}[1]{A_{*}(#1)}
\newcommand{\Ak}[2]{A^{#1}(#2)}
\newcommand{\Acycle}[2]{A_{#1}(#2)}

\newcommand{\chowring}[1]{A^{*}(#1)}
\newcommand{\cycleclass}[1]{[#1]}
\newcommand{\degree}{\operatorname{deg}}

% Stanley--Reisner and linear-relation ideals.
\newcommand{\ISR}{I_{\mathrm{SR}}}
\newcommand{\Ilin}{I_{\mathrm{lin}}}

% ==========================================================================
% Topology and compact torus notation
% ==========================================================================

\newcommand{\compacttorus}{(S^{1})^{n}}
\newcommand{\momentmap}{\mu}
\newcommand{\momentpolytope}{\Delta}

\newcommand{\piOne}[1]{\pi_{1}(#1)}
\newcommand{\Euler}[1]{\chi(#1)}

% ==========================================================================
% Morphisms and categorical notation
% ==========================================================================

\newcommand{\longto}{\longrightarrow}
\newcommand{\longfrom}{\longleftarrow}
\newcommand{\into}{\hookrightarrow}
\newcommand{\onto}{\twoheadrightarrow}
\newcommand{\isoto}{\xrightarrow{\sim}}
\newcommand{\birat}{\dashrightarrow}

\newcommand{\map}[3]{#1\colon #2\longrightarrow #3}
\newcommand{\rmap}[3]{#1\colon #2\dashrightarrow #3}

% Pullback and pushforward.
\newcommand{\pullback}[1]{#1^{*}}
\newcommand{\pushforward}[1]{#1_{*}}

% ==========================================================================
% Frequently used equations and constructions
% ==========================================================================

% The standard exact sequence for torus-invariant divisors:
%     M -> Z^{Sigma(1)} -> Cl(X_Sigma) -> 0.
\newcommand{\divisorsequence}{%
  M \longrightarrow \ZZ^{\Sigma(1)}
    \longrightarrow \Cl(X_{\Sigma}) \longrightarrow 0%
}

% Smooth toric divisor sequence:
\newcommand{\picardsequence}{%
  M \longrightarrow \ZZ^{\Sigma(1)}
    \longrightarrow \Pic(X_{\Sigma}) \longrightarrow 0%
}

% Character divisor formula.
\newcommand{\characterdivisorformula}{%
  \operatorname{div}(\chi^{m})
  =
  \sum_{\rho\in\Sigma(1)}
  \langle m,u_{\rho}\rangle D_{\rho}%
}

% Canonical divisor formula.
\newcommand{\canonicaldivisorformula}{%
  K_{X_{\Sigma}}
  =
  -\sum_{\rho\in\Sigma(1)}D_{\rho}%
}

% ==========================================================================
% Textual abbreviations
% ==========================================================================

\newcommand{\ie}{i.e.\ }
\newcommand{\eg}{e.g.\ }
\newcommand{\cf}{cf.\ }
\newcommand{\resp}{respectively}
\newcommand{\withoutloss}{without loss of generality}
\newcommand{\Withoutloss}{Without loss of generality}
\newcommand{\WLOG}{\text{without loss of generality}}

% ==========================================================================
% Optional notation aliases
% ==========================================================================
%
% These aliases are convenient but somewhat terse.  Keeping them grouped here
% makes it easy to remove them if a notation conflict appears later.

\newcommand{\Sig}{\Sigma}
\newcommand{\sig}{\sigma}
\newcommand{\taut}{\tau}
\newcommand{\rhoRay}{\rho}

% End of macros.tex


% ==========================================================================
% General conventions
% ==========================================================================

% Base field.  The default is an algebraically closed field k.
\providecommand{\Bbbk}{\mathbb{k}}

% Standard number systems.
\newcommand{\NN}{\mathbb{N}}
\newcommand{\ZZ}{\mathbb{Z}}
\newcommand{\QQ}{\mathbb{Q}}
\newcommand{\RR}{\mathbb{R}}
\newcommand{\CC}{\mathbb{C}}

% Affine and projective space.
\newcommand{\AAA}{\mathbb{A}}
\newcommand{\PP}{\mathbb{P}}

% Common script and calligraphic letters.
\newcommand{\cA}{\mathcal{A}}
\newcommand{\cB}{\mathcal{B}}
\newcommand{\cC}{\mathcal{C}}
\newcommand{\cD}{\mathcal{D}}
\newcommand{\cE}{\mathcal{E}}
\newcommand{\cF}{\mathcal{F}}
\newcommand{\cG}{\mathcal{G}}
\newcommand{\cH}{\mathcal{H}}
\newcommand{\cI}{\mathcal{I}}
\newcommand{\cJ}{\mathcal{J}}
\newcommand{\cK}{\mathcal{K}}
\newcommand{\cL}{\mathcal{L}}
\newcommand{\cM}{\mathcal{M}}
\newcommand{\cN}{\mathcal{N}}
\newcommand{\cO}{\mathcal{O}}
\newcommand{\cP}{\mathcal{P}}
\newcommand{\cQ}{\mathcal{Q}}
\newcommand{\cR}{\mathcal{R}}
\newcommand{\cS}{\mathcal{S}}
\newcommand{\cT}{\mathcal{T}}
\newcommand{\cU}{\mathcal{U}}
\newcommand{\cV}{\mathcal{V}}
\newcommand{\cW}{\mathcal{W}}
\newcommand{\cX}{\mathcal{X}}
\newcommand{\cY}{\mathcal{Y}}
\newcommand{\cZ}{\mathcal{Z}}

% Fraktur letters used occasionally for ideals and Lie algebras.
\newcommand{\fm}{\mathfrak{m}}
\newcommand{\fp}{\mathfrak{p}}
\newcommand{\fq}{\mathfrak{q}}
\newcommand{\fg}{\mathfrak{g}}
\newcommand{\ft}{\mathfrak{t}}

% Delimiters.
\DeclarePairedDelimiter{\paren}{(}{)}
\DeclarePairedDelimiter{\bracket}{[}{]}
\DeclarePairedDelimiter{\set}{\{}{\}}
\DeclarePairedDelimiter{\abs}{\lvert}{\rvert}
\DeclarePairedDelimiter{\norm}{\lVert}{\rVert}
\DeclarePairedDelimiter{\angles}{\langle}{\rangle}
\DeclarePairedDelimiter{\floor}{\lfloor}{\rfloor}
\DeclarePairedDelimiter{\ceil}{\lceil}{\rceil}

% Divisibility and restrictions.
\newcommand{\divides}{\mathrel{\mid}}
\newcommand{\ndivides}{\mathrel{\nmid}}
\newcommand{\restr}[2]{\left.#1\right|_{#2}}

% ==========================================================================
% Operators
% ==========================================================================

\DeclareMathOperator{\Spec}{Spec}
\DeclareMathOperator{\Proj}{Proj}
\DeclareMathOperator{\MaxSpec}{MaxSpec}
\DeclareMathOperator{\Hom}{Hom}
\DeclareMathOperator{\End}{End}
\DeclareMathOperator{\Aut}{Aut}
\DeclareMathOperator{\Isom}{Isom}
\DeclareMathOperator{\Mor}{Mor}
\DeclareMathOperator{\Ext}{Ext}
\DeclareMathOperator{\Tor}{Tor}
\DeclareMathOperator{\Pic}{Pic}
\DeclareMathOperator{\Cl}{Cl}
\DeclareMathOperator{\Div}{Div}
\DeclareMathOperator{\CaDiv}{CaDiv}
\DeclareMathOperator{\WDiv}{WDiv}
\DeclareMathOperator{\Prin}{Prin}
\DeclareMathOperator{\Supp}{Supp}
\DeclareMathOperator{\ord}{ord}
\DeclareMathOperator{\codim}{codim}
\DeclareMathOperator{\rank}{rank}
\DeclareMathOperator{\Span}{span}
\DeclareMathOperator{\cone}{Cone}
\DeclareMathOperator{\Conv}{Conv}
\DeclareMathOperator{\relint}{Relint}
\DeclareMathOperator{\Aff}{Aff}
\DeclareMathOperator{\vol}{Vol}
\DeclareMathOperator{\mult}{mult}
\DeclareMathOperator{\length}{length}
\DeclareMathOperator{\im}{im}
\DeclareMathOperator{\coker}{coker}
\DeclareMathOperator{\kerop}{ker}
\DeclareMathOperator{\id}{id}
\DeclareMathOperator{\diag}{diag}
\DeclareMathOperator{\Gal}{Gal}
\DeclareMathOperator{\Frac}{Frac}
\DeclareMathOperator{\trdeg}{trdeg}
\DeclareMathOperator{\Sing}{Sing}
\DeclareMathOperator{\Reg}{Reg}
\DeclareMathOperator{\Exc}{Exc}
\DeclareMathOperator{\Bl}{Bl}
\DeclareMathOperator{\NE}{NE}
\DeclareMathOperator{\Nef}{Nef}
\DeclareMathOperator{\Eff}{Eff}
\DeclareMathOperator{\Mov}{Mov}
\DeclareMathOperator{\Amp}{Amp}
\DeclareMathOperator{\Star}{Star}
\DeclareMathOperator{\Link}{Link}
\DeclareMathOperator{\Homog}{Homog}
\DeclareMathOperator{\Ann}{Ann}
\DeclareMathOperator{\ch}{ch}
\DeclareMathOperator{\td}{td}

% Prefer \ker for ordinary kernels; \kerop is supplied only when an operator
% with explicit limits is useful.
\renewcommand{\ker}{\operatorname{ker}}

% ==========================================================================
% Lattices, tori, characters, and one-parameter subgroups
% ==========================================================================

% CLS convention:
%   N is the lattice of one-parameter subgroups,
%   M = Hom_Z(N,Z) is the character lattice.
\newcommand{\Nlattice}{N}
\newcommand{\Mlattice}{M}

\newcommand{\NR}{N_{\RR}}
\newcommand{\NQ}{N_{\QQ}}
\newcommand{\MR}{M_{\RR}}
\newcommand{\MQ}{M_{\QQ}}

\newcommand{\NRbar}{\overline{N}_{\RR}}
\newcommand{\MRbar}{\overline{M}_{\RR}}

% Tensor extension of an arbitrary lattice.
\newcommand{\LR}[1]{#1_{\RR}}
\newcommand{\LQ}[1]{#1_{\QQ}}

% Dual lattice and dual vector space.
\newcommand{\dual}[1]{#1^{\vee}}
\newcommand{\ZZdual}[1]{\Hom_{\ZZ}(#1,\ZZ)}
\newcommand{\RRdual}[1]{\Hom_{\RR}(#1,\RR)}

% Pairing between M and N.
\newcommand{\pair}[2]{\left\langle #1,#2\right\rangle}

% Algebraic torus associated to N.
\newcommand{\TN}{T_N}
\newcommand{\TM}{T_M}
\newcommand{\torus}{T_N}
\newcommand{\Gm}{\mathbb{G}_{m}}
\newcommand{\Ga}{\mathbb{G}_{a}}

% Characters and one-parameter subgroups.
\newcommand{\character}[1]{\chi^{#1}}
\newcommand{\oneps}[1]{\lambda^{#1}}
\newcommand{\chiM}{\chi^m}
\newcommand{\lambdaU}{\lambda^u}

% Character and cocharacter groups.
\newcommand{\Xstar}[1]{X^{*}(#1)}
\newcommand{\Xsubstar}[1]{X_{*}(#1)}

% Group algebra / Laurent polynomial algebra.
\newcommand{\groupalg}[1]{\Bbbk[#1]}
\newcommand{\laurentring}[1]{\Bbbk[#1^{\pm 1}]}

% Weight spaces and gradings.
\newcommand{\weightspace}[2]{#1_{#2}}
\newcommand{\gradedpiece}[2]{#1_{#2}}
\newcommand{\wt}{\operatorname{wt}}

% ==========================================================================
% Semigroups, cones, and affine toric varieties
% ==========================================================================

% Semigroup generated by a subset A.
\newcommand{\semigroup}[1]{\NN #1}
\newcommand{\ZA}{\ZZ A}
\newcommand{\NA}{\NN A}

% Semigroup algebra.
\newcommand{\semigroupalg}[1]{\Bbbk[#1]}

% Dual cone and orthogonal complement.
\newcommand{\dualcone}[1]{#1^{\vee}}
\newcommand{\perpcone}[1]{#1^{\perp}}

% Cone attached to generators.
\newcommand{\Cone}[1]{\operatorname{Cone}(#1)}
\newcommand{\pos}[1]{\operatorname{pos}(#1)}

% Standard affine toric variety and semigroup associated to a cone.
\newcommand{\Ssigma}{S_{\sigma}}
\newcommand{\Usigma}{U_{\sigma}}
\newcommand{\Uface}[1]{U_{#1}}
\newcommand{\Xsigma}{X_{\sigma}}

% Generic versions.
\newcommand{\Scone}[1]{S_{#1}}
\newcommand{\Ucone}[1]{U_{#1}}
\newcommand{\Xcone}[1]{X_{#1}}

% Distinguished monomial.
\newcommand{\xmon}[1]{x^{#1}}

% Toric ideal associated to A.
\newcommand{\IA}{I_A}
\newcommand{\Iconfig}[1]{I_{#1}}

% Saturation.
\newcommand{\sat}[1]{\overline{#1}}
\newcommand{\saturation}[2]{#1:#2^{\infty}}

% Faces.
\newcommand{\faceleq}{\preceq}
\newcommand{\facet}{\text{facet}}

% ==========================================================================
% Fans and global toric varieties
% ==========================================================================

\newcommand{\fan}{\Sigma}
\newcommand{\fanp}{\Sigma'}
\newcommand{\fanq}{\Sigma''}

\newcommand{\Xfan}{X_{\Sigma}}
\newcommand{\Xfanp}{X_{\Sigma'}}
\newcommand{\Xof}[1]{X_{#1}}

% Cones of a fixed dimension.
\newcommand{\fancones}[1]{\Sigma(#1)}
\newcommand{\fanrays}{\Sigma(1)}
\newcommand{\fanmax}{\Sigma_{\max}}

% Support of a fan.
\newcommand{\fansupp}[1]{\lvert #1\rvert}

% Ray generator.  For rho in Sigma(1), u_rho denotes the primitive generator.
\newcommand{\urho}{u_{\rho}}
\newcommand{\uray}[1]{u_{#1}}
\newcommand{\vrho}{v_{\rho}}
\newcommand{\vray}[1]{v_{#1}}

% Orbit and orbit closure associated to a cone.
\newcommand{\Osigma}{O(\sigma)}
\newcommand{\Vsigma}{V(\sigma)}
\newcommand{\Ocone}[1]{O(#1)}
\newcommand{\Vcone}[1]{V(#1)}

% Star of a cone and quotient fan.
\newcommand{\starsigma}{\operatorname{Star}(\sigma)}
\newcommand{\starcone}[1]{\operatorname{Star}(#1)}

% Refinements and subdivisions.
\newcommand{\refines}{\succ}
\newcommand{\subdivision}[2]{#1 \to #2}

% ==========================================================================
% Polytopes and projective toric varieties
% ==========================================================================

\newcommand{\polytope}{P}
\newcommand{\polytopeQ}{Q}
\newcommand{\DeltaPoly}{\Delta}

\newcommand{\conv}[1]{\operatorname{Conv}(#1)}
\newcommand{\normalfan}[1]{\Sigma_{#1}}
\newcommand{\face}[2]{F_{#1}(#2)}

% Lattice points of a polytope.
\newcommand{\latticepoints}[1]{#1\cap M}
\newcommand{\numlatticepoints}[1]{\#(#1\cap M)}

% Polyhedron associated to a torus-invariant divisor.
\newcommand{\PD}{P_D}
\newcommand{\Pdiv}[1]{P_{#1}}

% Projective toric variety associated to a configuration or polytope.
\newcommand{\XA}{X_A}
\newcommand{\XP}{X_P}
\newcommand{\Xpoly}[1]{X_{#1}}

% Normal fan and support function.
\newcommand{\psiD}{\psi_D}
\newcommand{\supportfunction}[1]{\psi_{#1}}

% ==========================================================================
% Divisors, line bundles, and sections
% ==========================================================================

% Torus-invariant prime divisor associated to rho.
\newcommand{\Drho}{D_{\rho}}
\newcommand{\Dray}[1]{D_{#1}}

% Generic torus-invariant divisor.
\newcommand{\Dsum}{\sum_{\rho\in\Sigma(1)} a_{\rho}D_{\rho}}
\newcommand{\Dtoric}{D=\sum_{\rho\in\Sigma(1)}a_{\rho}D_{\rho}}

% Principal divisor of a character.
\newcommand{\divchi}[1]{\operatorname{div}(\chi^{#1})}

% Divisor and line-bundle notation.
\newcommand{\OX}{\mathcal{O}_X}
\newcommand{\OD}{\mathcal{O}_X(D)}
\newcommand{\Odiv}[1]{\mathcal{O}_X(#1)}
\newcommand{\Ldiv}[1]{\mathcal{L}(#1)}

% Global sections and cohomology.
\newcommand{\Hzero}[2]{H^{0}\!\left(#1,#2\right)}
\newcommand{\Hi}[3]{H^{#1}\!\left(#2,#3\right)}
\newcommand{\hzero}[2]{h^{0}\!\left(#1,#2\right)}
\newcommand{\hi}[3]{h^{#1}\!\left(#2,#3\right)}

% Canonical and anticanonical divisors.
\newcommand{\KX}{K_X}
\newcommand{\minusKX}{-K_X}
\newcommand{\omegaX}{\omega_X}

% Boundary divisor.
\newcommand{\boundaryX}{\partial X}
\newcommand{\toricboundary}{\sum_{\rho\in\Sigma(1)}D_{\rho}}

% Divisor class groups and numerical spaces.
\newcommand{\ClX}{\Cl(X)}
\newcommand{\PicX}{\Pic(X)}
\newcommand{\NSX}{\operatorname{NS}(X)}
\newcommand{\NOne}[1]{N^{1}(#1)_{\RR}}
\newcommand{\Ncurve}[1]{N_{1}(#1)_{\RR}}

% Section ring.
\newcommand{\sectionring}[2]{R(#1,#2)}
\newcommand{\RofD}{R(X,D)}

% ==========================================================================
% Cox rings and homogeneous coordinates
% ==========================================================================

\newcommand{\Cox}[1]{\operatorname{Cox}(#1)}
\newcommand{\CoxX}{\operatorname{Cox}(X_{\Sigma})}
\newcommand{\Scox}{S}
\newcommand{\Bfan}{B(\Sigma)}
\newcommand{\Zfan}{Z(\Sigma)}
\newcommand{\Uhat}{\widehat{X}}
\newcommand{\Gfan}{G_{\Sigma}}

% Cox variables.
\newcommand{\xrho}{x_{\rho}}
\newcommand{\xhat}[1]{\widehat{x}_{#1}}

% Homogeneous piece of degree alpha.
\newcommand{\Salpha}{S_{\alpha}}
\newcommand{\Sdegree}[1]{S_{#1}}

% ==========================================================================
% Toric surfaces and blowups
% ==========================================================================

% Hirzebruch surfaces.
\newcommand{\Hirz}[1]{\mathbb{F}_{#1}}
\newcommand{\Fn}{\mathbb{F}_{n}}

% Blowups.
\newcommand{\Blp}[2]{\operatorname{Bl}_{#2}(#1)}
\newcommand{\blowupmap}{\pi\colon \widetilde{X}\to X}
\newcommand{\Xtilde}{\widetilde{X}}
\newcommand{\Dtilde}{\widetilde{D}}
\newcommand{\Eexc}{E}

% Standard divisor classes on blowups of P^2.
\newcommand{\Hclass}{H}
\newcommand{\Eclass}[1]{E_{#1}}

% Self-intersection.
\newcommand{\selfint}[1]{#1^{2}}
\newcommand{\intersection}[2]{#1\cdot #2}

% Cyclically ordered primitive ray generators on a surface.
\newcommand{\vi}{v_i}
\newcommand{\vip}{v_{i+1}}
\newcommand{\vim}{v_{i-1}}

% ==========================================================================
% Singularities, resolutions, and birational geometry
% ==========================================================================

\newcommand{\resolution}{\pi\colon \widetilde{X}\to X}
\newcommand{\exceptional}{\operatorname{Exc}(\pi)}

\newcommand{\QCartier}{\QQ\text{-Cartier}}
\newcommand{\Qfactorial}{\QQ\text{-factorial}}
\newcommand{\QGorenstein}{\QQ\text{-Gorenstein}}

\newcommand{\MMP}{\mathrm{MMP}}
\newcommand{\NEbar}[1]{\overline{\operatorname{NE}}(#1)}
\newcommand{\Effbar}[1]{\overline{\operatorname{Eff}}(#1)}
\newcommand{\Movbar}[1]{\overline{\operatorname{Mov}}(#1)}

% Discrepancy notation.
\newcommand{\discrepancy}[2]{a(#1,#2)}

% ==========================================================================
% Intersection theory and Chow groups
% ==========================================================================

\newcommand{\Astar}[1]{A^{*}(#1)}
\newcommand{\Aunder}[1]{A_{*}(#1)}
\newcommand{\Ak}[2]{A^{#1}(#2)}
\newcommand{\Acycle}[2]{A_{#1}(#2)}

\newcommand{\chowring}[1]{A^{*}(#1)}
\newcommand{\cycleclass}[1]{[#1]}
\newcommand{\degree}{\operatorname{deg}}

% Stanley--Reisner and linear-relation ideals.
\newcommand{\ISR}{I_{\mathrm{SR}}}
\newcommand{\Ilin}{I_{\mathrm{lin}}}

% ==========================================================================
% Topology and compact torus notation
% ==========================================================================

\newcommand{\compacttorus}{(S^{1})^{n}}
\newcommand{\momentmap}{\mu}
\newcommand{\momentpolytope}{\Delta}

\newcommand{\piOne}[1]{\pi_{1}(#1)}
\newcommand{\Euler}[1]{\chi(#1)}

% ==========================================================================
% Morphisms and categorical notation
% ==========================================================================

\newcommand{\longto}{\longrightarrow}
\newcommand{\longfrom}{\longleftarrow}
\newcommand{\into}{\hookrightarrow}
\newcommand{\onto}{\twoheadrightarrow}
\newcommand{\isoto}{\xrightarrow{\sim}}
\newcommand{\birat}{\dashrightarrow}

\newcommand{\map}[3]{#1\colon #2\longrightarrow #3}
\newcommand{\rmap}[3]{#1\colon #2\dashrightarrow #3}

% Pullback and pushforward.
\newcommand{\pullback}[1]{#1^{*}}
\newcommand{\pushforward}[1]{#1_{*}}

% ==========================================================================
% Frequently used equations and constructions
% ==========================================================================

% The standard exact sequence for torus-invariant divisors:
%     M -> Z^{Sigma(1)} -> Cl(X_Sigma) -> 0.
\newcommand{\divisorsequence}{%
  M \longrightarrow \ZZ^{\Sigma(1)}
    \longrightarrow \Cl(X_{\Sigma}) \longrightarrow 0%
}

% Smooth toric divisor sequence:
\newcommand{\picardsequence}{%
  M \longrightarrow \ZZ^{\Sigma(1)}
    \longrightarrow \Pic(X_{\Sigma}) \longrightarrow 0%
}

% Character divisor formula.
\newcommand{\characterdivisorformula}{%
  \operatorname{div}(\chi^{m})
  =
  \sum_{\rho\in\Sigma(1)}
  \langle m,u_{\rho}\rangle D_{\rho}%
}

% Canonical divisor formula.
\newcommand{\canonicaldivisorformula}{%
  K_{X_{\Sigma}}
  =
  -\sum_{\rho\in\Sigma(1)}D_{\rho}%
}

% ==========================================================================
% Textual abbreviations
% ==========================================================================

\newcommand{\ie}{i.e.\ }
\newcommand{\eg}{e.g.\ }
\newcommand{\cf}{cf.\ }
\newcommand{\resp}{respectively}
\newcommand{\withoutloss}{without loss of generality}
\newcommand{\Withoutloss}{Without loss of generality}
\newcommand{\WLOG}{\text{without loss of generality}}

% ==========================================================================
% Optional notation aliases
% ==========================================================================
%
% These aliases are convenient but somewhat terse.  Keeping them grouped here
% makes it easy to remove them if a notation conflict appears later.

\newcommand{\Sig}{\Sigma}
\newcommand{\sig}{\sigma}
\newcommand{\taut}{\tau}
\newcommand{\rhoRay}{\rho}

% End of macros.tex


% ==========================================================================
% General conventions
% ==========================================================================

% Base field.  The default is an algebraically closed field k.
\providecommand{\Bbbk}{\mathbb{k}}

% Standard number systems.
\newcommand{\NN}{\mathbb{N}}
\newcommand{\ZZ}{\mathbb{Z}}
\newcommand{\QQ}{\mathbb{Q}}
\newcommand{\RR}{\mathbb{R}}
\newcommand{\CC}{\mathbb{C}}

% Affine and projective space.
\newcommand{\AAA}{\mathbb{A}}
\newcommand{\PP}{\mathbb{P}}

% Common script and calligraphic letters.
\newcommand{\cA}{\mathcal{A}}
\newcommand{\cB}{\mathcal{B}}
\newcommand{\cC}{\mathcal{C}}
\newcommand{\cD}{\mathcal{D}}
\newcommand{\cE}{\mathcal{E}}
\newcommand{\cF}{\mathcal{F}}
\newcommand{\cG}{\mathcal{G}}
\newcommand{\cH}{\mathcal{H}}
\newcommand{\cI}{\mathcal{I}}
\newcommand{\cJ}{\mathcal{J}}
\newcommand{\cK}{\mathcal{K}}
\newcommand{\cL}{\mathcal{L}}
\newcommand{\cM}{\mathcal{M}}
\newcommand{\cN}{\mathcal{N}}
\newcommand{\cO}{\mathcal{O}}
\newcommand{\cP}{\mathcal{P}}
\newcommand{\cQ}{\mathcal{Q}}
\newcommand{\cR}{\mathcal{R}}
\newcommand{\cS}{\mathcal{S}}
\newcommand{\cT}{\mathcal{T}}
\newcommand{\cU}{\mathcal{U}}
\newcommand{\cV}{\mathcal{V}}
\newcommand{\cW}{\mathcal{W}}
\newcommand{\cX}{\mathcal{X}}
\newcommand{\cY}{\mathcal{Y}}
\newcommand{\cZ}{\mathcal{Z}}

% Fraktur letters used occasionally for ideals and Lie algebras.
\newcommand{\fm}{\mathfrak{m}}
\newcommand{\fp}{\mathfrak{p}}
\newcommand{\fq}{\mathfrak{q}}
\newcommand{\fg}{\mathfrak{g}}
\newcommand{\ft}{\mathfrak{t}}

% Delimiters.
\DeclarePairedDelimiter{\paren}{(}{)}
\DeclarePairedDelimiter{\bracket}{[}{]}
\DeclarePairedDelimiter{\set}{\{}{\}}
\DeclarePairedDelimiter{\abs}{\lvert}{\rvert}
\DeclarePairedDelimiter{\norm}{\lVert}{\rVert}
\DeclarePairedDelimiter{\angles}{\langle}{\rangle}
\DeclarePairedDelimiter{\floor}{\lfloor}{\rfloor}
\DeclarePairedDelimiter{\ceil}{\lceil}{\rceil}

% Divisibility and restrictions.
\newcommand{\divides}{\mathrel{\mid}}
\newcommand{\ndivides}{\mathrel{\nmid}}
\newcommand{\restr}[2]{\left.#1\right|_{#2}}

% ==========================================================================
% Operators
% ==========================================================================

\DeclareMathOperator{\Spec}{Spec}
\DeclareMathOperator{\Proj}{Proj}
\DeclareMathOperator{\MaxSpec}{MaxSpec}
\DeclareMathOperator{\Hom}{Hom}
\DeclareMathOperator{\End}{End}
\DeclareMathOperator{\Aut}{Aut}
\DeclareMathOperator{\Isom}{Isom}
\DeclareMathOperator{\Mor}{Mor}
\DeclareMathOperator{\Ext}{Ext}
\DeclareMathOperator{\Tor}{Tor}
\DeclareMathOperator{\Pic}{Pic}
\DeclareMathOperator{\Cl}{Cl}
\DeclareMathOperator{\Div}{Div}
\DeclareMathOperator{\CaDiv}{CaDiv}
\DeclareMathOperator{\WDiv}{WDiv}
\DeclareMathOperator{\Prin}{Prin}
\DeclareMathOperator{\Supp}{Supp}
\DeclareMathOperator{\ord}{ord}
\DeclareMathOperator{\codim}{codim}
\DeclareMathOperator{\rank}{rank}
\DeclareMathOperator{\Span}{span}
\DeclareMathOperator{\cone}{Cone}
\DeclareMathOperator{\Conv}{Conv}
\DeclareMathOperator{\relint}{Relint}
\DeclareMathOperator{\Aff}{Aff}
\DeclareMathOperator{\vol}{Vol}
\DeclareMathOperator{\mult}{mult}
\DeclareMathOperator{\length}{length}
\DeclareMathOperator{\im}{im}
\DeclareMathOperator{\coker}{coker}
\DeclareMathOperator{\kerop}{ker}
\DeclareMathOperator{\id}{id}
\DeclareMathOperator{\diag}{diag}
\DeclareMathOperator{\Gal}{Gal}
\DeclareMathOperator{\Frac}{Frac}
\DeclareMathOperator{\trdeg}{trdeg}
\DeclareMathOperator{\Sing}{Sing}
\DeclareMathOperator{\Reg}{Reg}
\DeclareMathOperator{\Exc}{Exc}
\DeclareMathOperator{\Bl}{Bl}
\DeclareMathOperator{\NE}{NE}
\DeclareMathOperator{\Nef}{Nef}
\DeclareMathOperator{\Eff}{Eff}
\DeclareMathOperator{\Mov}{Mov}
\DeclareMathOperator{\Amp}{Amp}
\DeclareMathOperator{\Star}{Star}
\DeclareMathOperator{\Link}{Link}
\DeclareMathOperator{\Homog}{Homog}
\DeclareMathOperator{\Ann}{Ann}
\DeclareMathOperator{\ch}{ch}
\DeclareMathOperator{\td}{td}

% Prefer \ker for ordinary kernels; \kerop is supplied only when an operator
% with explicit limits is useful.
\renewcommand{\ker}{\operatorname{ker}}

% ==========================================================================
% Lattices, tori, characters, and one-parameter subgroups
% ==========================================================================

% CLS convention:
%   N is the lattice of one-parameter subgroups,
%   M = Hom_Z(N,Z) is the character lattice.
\newcommand{\Nlattice}{N}
\newcommand{\Mlattice}{M}

\newcommand{\NR}{N_{\RR}}
\newcommand{\NQ}{N_{\QQ}}
\newcommand{\MR}{M_{\RR}}
\newcommand{\MQ}{M_{\QQ}}

\newcommand{\NRbar}{\overline{N}_{\RR}}
\newcommand{\MRbar}{\overline{M}_{\RR}}

% Tensor extension of an arbitrary lattice.
\newcommand{\LR}[1]{#1_{\RR}}
\newcommand{\LQ}[1]{#1_{\QQ}}

% Dual lattice and dual vector space.
\newcommand{\dual}[1]{#1^{\vee}}
\newcommand{\ZZdual}[1]{\Hom_{\ZZ}(#1,\ZZ)}
\newcommand{\RRdual}[1]{\Hom_{\RR}(#1,\RR)}

% Pairing between M and N.
\newcommand{\pair}[2]{\left\langle #1,#2\right\rangle}

% Algebraic torus associated to N.
\newcommand{\TN}{T_N}
\newcommand{\TM}{T_M}
\newcommand{\torus}{T_N}
\newcommand{\Gm}{\mathbb{G}_{m}}
\newcommand{\Ga}{\mathbb{G}_{a}}

% Characters and one-parameter subgroups.
\newcommand{\character}[1]{\chi^{#1}}
\newcommand{\oneps}[1]{\lambda^{#1}}
\newcommand{\chiM}{\chi^m}
\newcommand{\lambdaU}{\lambda^u}

% Character and cocharacter groups.
\newcommand{\Xstar}[1]{X^{*}(#1)}
\newcommand{\Xsubstar}[1]{X_{*}(#1)}

% Group algebra / Laurent polynomial algebra.
\newcommand{\groupalg}[1]{\Bbbk[#1]}
\newcommand{\laurentring}[1]{\Bbbk[#1^{\pm 1}]}

% Weight spaces and gradings.
\newcommand{\weightspace}[2]{#1_{#2}}
\newcommand{\gradedpiece}[2]{#1_{#2}}
\newcommand{\wt}{\operatorname{wt}}

% ==========================================================================
% Semigroups, cones, and affine toric varieties
% ==========================================================================

% Semigroup generated by a subset A.
\newcommand{\semigroup}[1]{\NN #1}
\newcommand{\ZA}{\ZZ A}
\newcommand{\NA}{\NN A}

% Semigroup algebra.
\newcommand{\semigroupalg}[1]{\Bbbk[#1]}

% Dual cone and orthogonal complement.
\newcommand{\dualcone}[1]{#1^{\vee}}
\newcommand{\perpcone}[1]{#1^{\perp}}

% Cone attached to generators.
\newcommand{\Cone}[1]{\operatorname{Cone}(#1)}
\newcommand{\pos}[1]{\operatorname{pos}(#1)}

% Standard affine toric variety and semigroup associated to a cone.
\newcommand{\Ssigma}{S_{\sigma}}
\newcommand{\Usigma}{U_{\sigma}}
\newcommand{\Uface}[1]{U_{#1}}
\newcommand{\Xsigma}{X_{\sigma}}

% Generic versions.
\newcommand{\Scone}[1]{S_{#1}}
\newcommand{\Ucone}[1]{U_{#1}}
\newcommand{\Xcone}[1]{X_{#1}}

% Distinguished monomial.
\newcommand{\xmon}[1]{x^{#1}}

% Toric ideal associated to A.
\newcommand{\IA}{I_A}
\newcommand{\Iconfig}[1]{I_{#1}}

% Saturation.
\newcommand{\sat}[1]{\overline{#1}}
\newcommand{\saturation}[2]{#1:#2^{\infty}}

% Faces.
\newcommand{\faceleq}{\preceq}
\newcommand{\facet}{\text{facet}}

% ==========================================================================
% Fans and global toric varieties
% ==========================================================================

\newcommand{\fan}{\Sigma}
\newcommand{\fanp}{\Sigma'}
\newcommand{\fanq}{\Sigma''}

\newcommand{\Xfan}{X_{\Sigma}}
\newcommand{\Xfanp}{X_{\Sigma'}}
\newcommand{\Xof}[1]{X_{#1}}

% Cones of a fixed dimension.
\newcommand{\fancones}[1]{\Sigma(#1)}
\newcommand{\fanrays}{\Sigma(1)}
\newcommand{\fanmax}{\Sigma_{\max}}

% Support of a fan.
\newcommand{\fansupp}[1]{\lvert #1\rvert}

% Ray generator.  For rho in Sigma(1), u_rho denotes the primitive generator.
\newcommand{\urho}{u_{\rho}}
\newcommand{\uray}[1]{u_{#1}}
\newcommand{\vrho}{v_{\rho}}
\newcommand{\vray}[1]{v_{#1}}

% Orbit and orbit closure associated to a cone.
\newcommand{\Osigma}{O(\sigma)}
\newcommand{\Vsigma}{V(\sigma)}
\newcommand{\Ocone}[1]{O(#1)}
\newcommand{\Vcone}[1]{V(#1)}

% Star of a cone and quotient fan.
\newcommand{\starsigma}{\operatorname{Star}(\sigma)}
\newcommand{\starcone}[1]{\operatorname{Star}(#1)}

% Refinements and subdivisions.
\newcommand{\refines}{\succ}
\newcommand{\subdivision}[2]{#1 \to #2}

% ==========================================================================
% Polytopes and projective toric varieties
% ==========================================================================

\newcommand{\polytope}{P}
\newcommand{\polytopeQ}{Q}
\newcommand{\DeltaPoly}{\Delta}

\newcommand{\conv}[1]{\operatorname{Conv}(#1)}
\newcommand{\normalfan}[1]{\Sigma_{#1}}
\newcommand{\face}[2]{F_{#1}(#2)}

% Lattice points of a polytope.
\newcommand{\latticepoints}[1]{#1\cap M}
\newcommand{\numlatticepoints}[1]{\#(#1\cap M)}

% Polyhedron associated to a torus-invariant divisor.
\newcommand{\PD}{P_D}
\newcommand{\Pdiv}[1]{P_{#1}}

% Projective toric variety associated to a configuration or polytope.
\newcommand{\XA}{X_A}
\newcommand{\XP}{X_P}
\newcommand{\Xpoly}[1]{X_{#1}}

% Normal fan and support function.
\newcommand{\psiD}{\psi_D}
\newcommand{\supportfunction}[1]{\psi_{#1}}

% ==========================================================================
% Divisors, line bundles, and sections
% ==========================================================================

% Torus-invariant prime divisor associated to rho.
\newcommand{\Drho}{D_{\rho}}
\newcommand{\Dray}[1]{D_{#1}}

% Generic torus-invariant divisor.
\newcommand{\Dsum}{\sum_{\rho\in\Sigma(1)} a_{\rho}D_{\rho}}
\newcommand{\Dtoric}{D=\sum_{\rho\in\Sigma(1)}a_{\rho}D_{\rho}}

% Principal divisor of a character.
\newcommand{\divchi}[1]{\operatorname{div}(\chi^{#1})}

% Divisor and line-bundle notation.
\newcommand{\OX}{\mathcal{O}_X}
\newcommand{\OD}{\mathcal{O}_X(D)}
\newcommand{\Odiv}[1]{\mathcal{O}_X(#1)}
\newcommand{\Ldiv}[1]{\mathcal{L}(#1)}

% Global sections and cohomology.
\newcommand{\Hzero}[2]{H^{0}\!\left(#1,#2\right)}
\newcommand{\Hi}[3]{H^{#1}\!\left(#2,#3\right)}
\newcommand{\hzero}[2]{h^{0}\!\left(#1,#2\right)}
\newcommand{\hi}[3]{h^{#1}\!\left(#2,#3\right)}

% Canonical and anticanonical divisors.
\newcommand{\KX}{K_X}
\newcommand{\minusKX}{-K_X}
\newcommand{\omegaX}{\omega_X}

% Boundary divisor.
\newcommand{\boundaryX}{\partial X}
\newcommand{\toricboundary}{\sum_{\rho\in\Sigma(1)}D_{\rho}}

% Divisor class groups and numerical spaces.
\newcommand{\ClX}{\Cl(X)}
\newcommand{\PicX}{\Pic(X)}
\newcommand{\NSX}{\operatorname{NS}(X)}
\newcommand{\NOne}[1]{N^{1}(#1)_{\RR}}
\newcommand{\Ncurve}[1]{N_{1}(#1)_{\RR}}

% Section ring.
\newcommand{\sectionring}[2]{R(#1,#2)}
\newcommand{\RofD}{R(X,D)}

% ==========================================================================
% Cox rings and homogeneous coordinates
% ==========================================================================

\newcommand{\Cox}[1]{\operatorname{Cox}(#1)}
\newcommand{\CoxX}{\operatorname{Cox}(X_{\Sigma})}
\newcommand{\Scox}{S}
\newcommand{\Bfan}{B(\Sigma)}
\newcommand{\Zfan}{Z(\Sigma)}
\newcommand{\Uhat}{\widehat{X}}
\newcommand{\Gfan}{G_{\Sigma}}

% Cox variables.
\newcommand{\xrho}{x_{\rho}}
\newcommand{\xhat}[1]{\widehat{x}_{#1}}

% Homogeneous piece of degree alpha.
\newcommand{\Salpha}{S_{\alpha}}
\newcommand{\Sdegree}[1]{S_{#1}}

% ==========================================================================
% Toric surfaces and blowups
% ==========================================================================

% Hirzebruch surfaces.
\newcommand{\Hirz}[1]{\mathbb{F}_{#1}}
\newcommand{\Fn}{\mathbb{F}_{n}}

% Blowups.
\newcommand{\Blp}[2]{\operatorname{Bl}_{#2}(#1)}
\newcommand{\blowupmap}{\pi\colon \widetilde{X}\to X}
\newcommand{\Xtilde}{\widetilde{X}}
\newcommand{\Dtilde}{\widetilde{D}}
\newcommand{\Eexc}{E}

% Standard divisor classes on blowups of P^2.
\newcommand{\Hclass}{H}
\newcommand{\Eclass}[1]{E_{#1}}

% Self-intersection.
\newcommand{\selfint}[1]{#1^{2}}
\newcommand{\intersection}[2]{#1\cdot #2}

% Cyclically ordered primitive ray generators on a surface.
\newcommand{\vi}{v_i}
\newcommand{\vip}{v_{i+1}}
\newcommand{\vim}{v_{i-1}}

% ==========================================================================
% Singularities, resolutions, and birational geometry
% ==========================================================================

\newcommand{\resolution}{\pi\colon \widetilde{X}\to X}
\newcommand{\exceptional}{\operatorname{Exc}(\pi)}

\newcommand{\QCartier}{\QQ\text{-Cartier}}
\newcommand{\Qfactorial}{\QQ\text{-factorial}}
\newcommand{\QGorenstein}{\QQ\text{-Gorenstein}}

\newcommand{\MMP}{\mathrm{MMP}}
\newcommand{\NEbar}[1]{\overline{\operatorname{NE}}(#1)}
\newcommand{\Effbar}[1]{\overline{\operatorname{Eff}}(#1)}
\newcommand{\Movbar}[1]{\overline{\operatorname{Mov}}(#1)}

% Discrepancy notation.
\newcommand{\discrepancy}[2]{a(#1,#2)}

% ==========================================================================
% Intersection theory and Chow groups
% ==========================================================================

\newcommand{\Astar}[1]{A^{*}(#1)}
\newcommand{\Aunder}[1]{A_{*}(#1)}
\newcommand{\Ak}[2]{A^{#1}(#2)}
\newcommand{\Acycle}[2]{A_{#1}(#2)}

\newcommand{\chowring}[1]{A^{*}(#1)}
\newcommand{\cycleclass}[1]{[#1]}
\newcommand{\degree}{\operatorname{deg}}

% Stanley--Reisner and linear-relation ideals.
\newcommand{\ISR}{I_{\mathrm{SR}}}
\newcommand{\Ilin}{I_{\mathrm{lin}}}

% ==========================================================================
% Topology and compact torus notation
% ==========================================================================

\newcommand{\compacttorus}{(S^{1})^{n}}
\newcommand{\momentmap}{\mu}
\newcommand{\momentpolytope}{\Delta}

\newcommand{\piOne}[1]{\pi_{1}(#1)}
\newcommand{\Euler}[1]{\chi(#1)}

% ==========================================================================
% Morphisms and categorical notation
% ==========================================================================

\newcommand{\longto}{\longrightarrow}
\newcommand{\longfrom}{\longleftarrow}
\newcommand{\into}{\hookrightarrow}
\newcommand{\onto}{\twoheadrightarrow}
\newcommand{\isoto}{\xrightarrow{\sim}}
\newcommand{\birat}{\dashrightarrow}

\newcommand{\map}[3]{#1\colon #2\longrightarrow #3}
\newcommand{\rmap}[3]{#1\colon #2\dashrightarrow #3}

% Pullback and pushforward.
\newcommand{\pullback}[1]{#1^{*}}
\newcommand{\pushforward}[1]{#1_{*}}

% ==========================================================================
% Frequently used equations and constructions
% ==========================================================================

% The standard exact sequence for torus-invariant divisors:
%     M -> Z^{Sigma(1)} -> Cl(X_Sigma) -> 0.
\newcommand{\divisorsequence}{%
  M \longrightarrow \ZZ^{\Sigma(1)}
    \longrightarrow \Cl(X_{\Sigma}) \longrightarrow 0%
}

% Smooth toric divisor sequence:
\newcommand{\picardsequence}{%
  M \longrightarrow \ZZ^{\Sigma(1)}
    \longrightarrow \Pic(X_{\Sigma}) \longrightarrow 0%
}

% Character divisor formula.
\newcommand{\characterdivisorformula}{%
  \operatorname{div}(\chi^{m})
  =
  \sum_{\rho\in\Sigma(1)}
  \langle m,u_{\rho}\rangle D_{\rho}%
}

% Canonical divisor formula.
\newcommand{\canonicaldivisorformula}{%
  K_{X_{\Sigma}}
  =
  -\sum_{\rho\in\Sigma(1)}D_{\rho}%
}

% ==========================================================================
% Textual abbreviations
% ==========================================================================

\newcommand{\ie}{i.e.\ }
\newcommand{\eg}{e.g.\ }
\newcommand{\cf}{cf.\ }
\newcommand{\resp}{respectively}
\newcommand{\withoutloss}{without loss of generality}
\newcommand{\Withoutloss}{Without loss of generality}
\newcommand{\WLOG}{\text{without loss of generality}}

% ==========================================================================
% Optional notation aliases
% ==========================================================================
%
% These aliases are convenient but somewhat terse.  Keeping them grouped here
% makes it easy to remove them if a notation conflict appears later.

\newcommand{\Sig}{\Sigma}
\newcommand{\sig}{\sigma}
\newcommand{\taut}{\tau}
\newcommand{\rhoRay}{\rho}

% End of macros.tex
